\chapter{Parametrizaciónes básicas}

En este capítulo se presentan algunos ejemplos sencillos sobre la manera de representar curvas,
superficies o sólidos cuyas generalidades se presentaron en la sección \secref{regionesparametrizadas}. \index{Parametrización}
Con las ecuaciones (\ref{param1}) y (\ref{param2}) que corresponden a ideas muy simples, se construirán otras parametrizaciones y luego, con ayuda de estas
funciones, se parametrizarán regiones en el espacio. Se hacen énfasis en la construcción de este 
tipo de funciones, pues en los capítulos posteriores se hará uso de ellas sin profundizar demasiado en su obtención. 

\section{Curvas} \index{Parametrización!de curvas} \index{Curva} 
De acuerdo con \cite{marsden_calculo_2013}, una curva $C$ en $\mathbb{R}^n$ es la imagen de una función
contínua a trozos, $\ve{r}:[a,b]\rightarrow \mathbb{R}^n$. 
A la función $\ve{r}$ la denominaremos \emph{una parametrización} de la curva $C$. 
\vskip0.3cm 
Si $g:[c,d]\rightarrow \lbrack a,b]$ es una función biyectiva y continua en $[c,d]$
tal que $\alpha =r\circ g, \,$ a $\alpha,$ se le denomina una \emph{reparametrización} de $\ve{r}.$
Si la función $g$ es creciente $\left(g^{\prime }(t)>0\right), $
entonces $\alpha $ preserva la orientación, y si $g$ es decreciente
$\left( g^{\prime }(t)<0\right), $ entonces $\alpha $ invierte la orientación de la curva.
\vskip0.1cm
Dada una curva $C$ parametrizada por medio de la función $\ve{r}(t)$, con $t \in [a,b],$
la sustitución 
\begin{equation} \label{param3}
\ve{v}(t)=\ve{r}(a+b-t)
\end{equation}
parametriza la misma curva con sentido contrario para $t \in [a,b].$ \index{Cambio de sentido!Curvas}
%\vskip0.1cm

\vskip0.2cm
\noindent
\begin{minipage}{\textwidth}
	\centering
	\captionof{figure}{${\bf u}(t)=\langle t^3-3t, t \rangle$}
	\begin{minipage}{0.35\linewidth}
		\centering
		\includegraphics[width=\linewidth]{Fig/20.pdf}
	\end{minipage}
	
	\vspace{0.3em}
	\parbox{0.95\linewidth}{\centering\fontsize{10pt}{10pt}\selectfont\textbf{Fuente:} elaboración propia.}
\end{minipage}
\vskip0.2cm
 
 
 
La curva determinada por la ecuación $x=y^3-3y$, con  $y\in[-1,2]$
se puede parametrizar mediante la función 
$\ve{u}(t)=\langle t^3-3t, t \rangle$
con $t \in [-1,2]$. El punto inicial 
es  $(2,-1)$ y el punto final es $(2,2).$

Al cambiar $t$ por $-1+2-t=1-t$ y utilizar la expresión (\ref{param3}), 
se establece que la función $\ve{v}(t)= \la -t^3+3t^2-2,1-t\ra$
con $t \in [-1,2]$, que representa la misma curva, recorrida en sentido contrario.

\vskip0.2cm
\noindent
\begin{minipage}{\textwidth}
	\centering
	\captionof{figure}{${\bf u}(t)=\langle  t,t^4-3t^2 \rangle$}
	\begin{minipage}{0.35\linewidth}
		\centering
		\includegraphics[width=\linewidth]{Fig/21.pdf}
	\end{minipage}
	
	\vspace{0.3em}
	\parbox{0.95\linewidth}{\centering\fontsize{10pt}{10pt}\selectfont\textbf{Fuente:} elaboración propia.}
\end{minipage}
\vskip0.2cm

%\vskip0.1cm 
La función $y= x^4-3x^2,$ con $x\in[-2,2]$
se parametriza mediante la expresión 
$\ve{r}(t)=\langle t,t^4-3t^2 \rangle$ con $t \in [-2,2]$; el punto inicial de la curva descrita 
es  $(-2,4)$ y el punto final es $(2,4).$
Esta representación no es única, pues $\ve{u}(t)=\langle 2t,16t^4-12t^2 \rangle$ 
con $t \in [-1,1]$ representa la misma curva, con el mismo recorrido y los puntos terminales.

\subsection*{Tres ideas elementales}
Las siguientes tres parametrizaciones se usarán en gran parte de los problemas expuestos y servirán
como insumo para construir o representar varias regiones en el espacio.

\begin{tcolorbox}[breakable,enhanced,title=\bf Parametrización de un segmento]
Para dos puntos $a$ y $b$ en $\mathbb{R}^{n},$  la función 
\begin{equation} \label{param1}
\overrightarrow{r}(t)=(1-t) a+tb,\,\, t\in \lbrack 0,1]
\end{equation}
parametriza el segmento rectilíneo que une a $a$ con $b$. 
\end{tcolorbox}
Intercambiando $a$ y $b$ en la fórmula (\ref{param1}) se logra cambiar el punto inicial con el punto final.
Una expresión como esta la denominaremos una 
\emph{combinación convexa} entre los puntos $a$ y $b$. \index{Combinación convexa}

\begin{tcolorbox}[breakable,enhanced,title=\bf Parametrización de la circunferencia unitaria]
La circunferencia unitaria cuya ecuación es $x^2+y^2=1,$
puede parametrizarse en la forma 
\begin{equation} \label{param2}
\ve{r}(t)=\left\langle \cos(t), \sin(t)\right\rangle,\, t\in [0,2\pi].
\end{equation}
La identidad $\cos^2t + \sin^2 t =1$ motiva el uso de esta elección.
\end{tcolorbox}
En la expresión (\ref{param2}) también se puede asignar $x=\sin (t) $ y $y=\cos (t) $, de esta manera se describe la misma curva; la diferencia está en los puntos terminales o en la orientación.


\vskip0.2cm
\noindent
\begin{minipage}{\textwidth}
	\centering
	\captionof{figure}{Gráficos de ${\bf u}(t)=\langle\cos(t),\sin(t)\rangle$ y ${\bf u}(t)=\langle\sin(t),\cos(t)\rangle$}
	\begin{minipage}{0.40\linewidth}
		\centering
		\vspace{-2mm}
		\includegraphics[width=\linewidth]{Fig/22.pdf}
	\end{minipage}
	
	\vspace{0.3em}
	\parbox{0.95\linewidth}{\centering\fontsize{9pt}{10pt}\selectfont\textbf{Fuente:} elaboración propia.}
\end{minipage}
\vskip0.2cm


Para una misma curva, la parametrización presentada podría no ser única.
Por ejemplo, en el caso de la circunferencia unitaria también se puede \mbox{representar} por medio de la función
${\bf u}(t)=\langle \cos(3t), \sin(3t)\rangle$, $t \in [0,2\pi/3]$.

%\vspace{-10mm}

\begin{tcolorbox}[breakable,enhanced,title=\bf Parametrización de la hipérbola unitaria]
La ecuación $x^2-y^2=1$ representa una hipérbola cuyas ramas abren hacia el eje $x$.
Una porción de esta curva se parametriza por medio de la función 
$\ve{r}(t)=\la\cosh(t), \sinh(t)\ra,$ con $t\in [-1,1]$. La definición de esta función es motivada
por la relación $\cosh^2 t - \sinh^2 t =1.$ 
\end{tcolorbox}

%\vspace{-2mm}

La otra rama se representa por la función 
{\small $\ve{r}(t)=\la -\cosh(t), \sinh(t)\ra.$}  
Las siguientes parametrizaciones representan solo una rama de la hipérbola.
\vskip0.2cm
Con los elementos expuestos presentaremos otras curvas con su respectiva parametrización.

\vskip0.2cm
\noindent
\begin{minipage}{\textwidth}
	\centering
	\captionof{figure}{Gráficos de ${\bf r}(t)=\la \cosh(t),\sinh(t)\ra$ y ${\bf r}(t)=\la \sinh(t),\cosh(t)\ra$}
	\begin{minipage}{0.40\linewidth}
		\centering
		\includegraphics[width=\linewidth]{Fig/23.pdf}
	\end{minipage}
	
	\vspace{0.3em}
	\parbox{0.95\linewidth}{\centering\fontsize{9pt}{10pt}\selectfont\textbf{Fuente:} elaboración propia.}
\end{minipage}
\vskip0.2cm


\subsection*{Ejemplos de parametrización de curvas}

\vspace{0pt}
\textcolor{black}{\scriptsize\ding{110}}\quad El segmento de recta que une los puntos $A(1,1)$ y $B(3,2)$, iniciando en $A$ y terminando en $B$, se parametriza al simplificar la combinación convexa
{\small
	\setlength{\abovedisplayskip}{2pt}
	\setlength{\belowdisplayskip}{2pt}
	\[ (1-t)(1,1) + t(3,2), \]
}
la cual define la función
{\small
	\setlength{\abovedisplayskip}{2pt}
	\[ \ve{r}(t) = \left\langle 1+2t,1+t\right\rangle, \quad t \in [0,1]. \]
}
	
%\vskip0.1cm
\noindent
\begin{minipage}{\textwidth}
	\centering
	\captionof{figure}{Segmento en el plano}
	\begin{minipage}{0.24\linewidth}
		\centering
		\includegraphics[width=\linewidth]{Fig/24.pdf}
	\end{minipage}
	
	\vspace{0.3em}
	\parbox{0.95\linewidth}{\centering\fontsize{9pt}{10pt}\selectfont\textbf{Fuente:} elaboración propia.}
\end{minipage}
\vskip0.2cm		


\vspace{0pt}
\textcolor{black}{\scriptsize\ding{110}}\quad El segmento rectilíneo con punto inicial {\small $(2,-1,0)$} y punto final {\small $(-1, 2, 2),$} se parametriza simplificando una combinación convexa de los puntos
\[
(1 - t)(2,-1,0) + t(-1,2,2),
\]
lo cual permite definir la función
\[
\ve{r}(t) = \left\langle 2 - 2t,\,-1 + 3t,\, 2t \right\rangle, \quad t \in [0,1].
\]

\vskip0.2cm
\noindent
\begin{minipage}{\textwidth}
	\centering
	\captionof{figure}{Segmento en el espacio}
	\label{segmento_espacio}
	\begin{minipage}{0.30\linewidth}
		\centering
		\includegraphics[width=\linewidth]{Fig/18.pdf}
	\end{minipage}
	
	\vspace{0.3em}
	\parbox{0.95\linewidth}{\centering\fontsize{8pt}{10pt}\selectfont\textbf{Fuente:} elaboración propia.}
\end{minipage}
\vskip0.2cm




\textcolor{black}{\scriptsize\ding{110}}\quad La elipse centrada en el origen y alargada en el sentido de las $x$  
con ecuación $\frac{x^2}{9} + \frac{y^2}{4}=1$ se parametriza al hacer $x=3 \cos(t)$ y $y=2 \sin(t)$.  
Es decir, esta curva se representa por medio de la función  
$$\ve{r}(t)=\left\langle 3 \cos(t), 2\sin(t)\right\rangle,\ t\in [0,2\pi].$$

\vskip0.2cm
\noindent
\begin{minipage}{\textwidth}
	\centering
	\captionof{figure}{${\bf r}(t)=\langle 3\cos(t),2\sin(t)\rangle$}
	\begin{minipage}{0.28\linewidth}
		\centering
		\includegraphics[width=\linewidth]{Fig/25.pdf}
	\end{minipage}
	
	\vspace{0.3em}
	\parbox{0.95\linewidth}{\centering\fontsize{9pt}{10pt}\selectfont\textbf{Fuente:} elaboración propia.}
\end{minipage}
\vskip0.2cm


\vspace{0pt}
\textcolor{black}{\scriptsize\ding{110}}\quad Con $t \in [0,\pi]$ se describe la mitad de la elipse. Utilizando la fórmula (\ref{param3}), se obtiene $x = -3\cos(t)$ y $y = 2\sin(t)$. De esta manera se describe la misma curva, variando sus puntos terminales y la orientación.

\vskip0.2cm
\noindent
\begin{minipage}{\textwidth}
	\centering
	\captionof{figure}{${\bf r}(t) = \langle -3\cos(t),\, 2\sin(t) \rangle$}
	\label{semielipse_izquierda}
	\begin{minipage}{0.30\linewidth}
		\centering
		\includegraphics[width=\linewidth]{Fig/26.pdf}
	\end{minipage}
	
	\vspace{0.3em}
	\parbox{0.95\linewidth}{\centering\fontsize{8pt}{10pt}\selectfont\textbf{Fuente:} elaboración propia.}
\end{minipage}
\vskip0.2cm


\vspace{0pt}
\textcolor{black}{\scriptsize\ding{110}}\quad En el espacio, la circunferencia de radio $2$ centrada
en el punto $(0, 2, 0)$ y paralela al plano $xz,$
se puede parametrizar como 
$\ve{r}(t) = \la 2\cos t, 2, 2\sin t\ra,$ $t\in[0,2\pi].$ 
Obsérvese que esta curva se halla en el plano $y=2$.

\vskip0.2cm
\noindent
\begin{minipage}{\textwidth}
	\centering
	\captionof{figure}{Circunferencia en el espacio}
	\label{Circunferencia_espacio}
	\begin{minipage}{0.30\linewidth}
		\centering
		\includegraphics[width=\linewidth]{Fig/19.pdf}
	\end{minipage}
	
	\vspace{0.3em}
	\parbox{0.95\linewidth}{\centering\fontsize{8pt}{10pt}\selectfont\textbf{Fuente:} elaboración propia.}
\end{minipage}
\vskip0.2cm


\vspace{0pt}
\textcolor{black}{\scriptsize\ding{110}}\quad Las ecuaciones $x-y+3z=3$ y $x^2+y^2=4$ representan un plano y un cilindro. En el plano $xy$ la ecuación del cilindro corresponde a una circunferencia de radio $2$ que se parametriza con $x=2 \cos t,$ $y=2 \sin t.$ De la ecuación del plano se obtiene $z$.
Luego, la intersección de las superficies está dada por $\ve{r}(t)=\la 2\cos t, 2\sin t,1-\frac{2}{3}(\cos t+\sin t) \ra.$


\vskip0.2cm
\noindent
\begin{minipage}{\textwidth}
	\centering
	\captionof{figure}{Elipse en el espacio}
	\label{Elipse_espacio}
	\begin{minipage}{0.28\linewidth}
		\centering
		\includegraphics[width=\linewidth]{Fig/27.pdf}
	\end{minipage}
	
	\vspace{0.3em}
	\parbox{0.95\linewidth}{\centering\fontsize{8pt}{10pt}\selectfont\textbf{Fuente:} elaboración propia.}
\end{minipage}
\vskip0.2cm


\subsection*{Uniendo dos puntos con una curva}
Para dos puntos $P$ y $Q$ en $\mathbb{R}^n$, si $f$ y $g$ son dos funciones contínuas en $[a,b]$ tales que 
$f(a)=1,$ $g(a)=0$ y $f(b)=0,$ $g(b)=1,$ entonces la función
\begin{equation}
\ve{r}(t)=f(t)P+g(t)Q,\,\, t \in[a,b].
\end{equation}
parametriza la curva que une a $P$ con $Q$. 
%Además $\ve{r}(t)$ satisface que $\ve{r}(a)=P$ y $ \ve{r}(b)=Q;$
La forma de la trayectoria depende de las funciones $f$ y $g.$
Enseguida ilustraremos lo anterior. 


\vspace{0pt}
\textcolor{black}{\scriptsize\ding{110}}\quad Con las funciones $g(t)=\sqrt[4]{t}$ 
y $f(t)=(1-t)^3$, uniremos los puntos $A(1,1)$ y $B(3,2)$ mediante la expresión 
$(1-t)^3 A + \sqrt[4]{t} B.$ 
\vskip0.05cm
Esto permite definir la función 
${\bf r}(t)=(1-t) ^{3}(1,1) +\sqrt[4]{t}(3,2),$
con  $t\in[0,1].$


\vskip0.2cm
\noindent
\begin{minipage}{\textwidth}
	\centering
	\captionof{figure}{${\bf r}(t)=(1-t) ^{3}(1,1) +\sqrt[4]{t}(3,2)$}
	\begin{minipage}{0.26\linewidth}
		\centering
		\includegraphics[width=\linewidth]{Fig/28.pdf}
	\end{minipage}
	
	\vspace{0.3em}
	\parbox{0.95\linewidth}{\centering\fontsize{8pt}{10pt}\selectfont\textbf{Fuente:} elaboración propia.}
\end{minipage}
\vskip0.2cm


\vspace{0pt}
\textcolor{black}{\scriptsize\ding{110}}\quad Con ayuda de las funciones $g(t) = t^6$ y $f(t) = (1 - t^2)^3$, se define la función
\[
\ve{r}(t) = (1 - t^2)^3(1,1) + t^6(3,2),
\]
que une los puntos $A(1,1)$ y $B(3,2)$ mediante la combinación
\[
(1 - t^2)^3 A + t^6 B, \quad t \in [0,1].
\]


\vskip0.2cm
\noindent
\begin{minipage}{\textwidth}
	\centering
	\captionof{figure}{${\bf r}(t) = (1 - t^2)^3(1,1) + t^6(3,2)$}
	\begin{minipage}{0.26\linewidth}
		\centering
		\includegraphics[width=\linewidth]{Fig/29.pdf}
	\end{minipage}
	
	\vspace{0.3em}
	\parbox{0.95\linewidth}{\centering\fontsize{8pt}{10pt}\selectfont\textbf{Fuente:} elaboración propia.}
\end{minipage}
\vskip0.2cm


\vspace{0pt}
\textcolor{black}{\scriptsize\ding{110}}\quad Utilizando una interpolación cuadrática entre los puntos $(1,1),$ $(3,2)$ y un punto intermedio 
entre ellos, se obtienen las siguientes funciones: \\[0.2cm]
$y=\frac{51-61t+16t^2}{6}$, $y=\frac{27-29t+8t^2}{6}$, $y=\frac{22-27t+7t^2}{6}$, 
$y=\frac{-16+33t-7t^2}{10},$ $y=\frac{-5+16t-3t^2}{8}$. \\[0.2cm]
\\Se muestran estas curvas parametrizadas con
$\la t, y(t)\ra $, $t\in[1,3].$

\vskip0.2cm
\noindent
\begin{minipage}{\textwidth}
	\centering
	\captionof{figure}{Cinco polinomios que unen $A$ con $B$}
	\begin{minipage}{0.29\linewidth}
		\centering
		\includegraphics[width=\linewidth]{Fig/30.pdf}
	\end{minipage}
	
	\vspace{0.3em}
	\parbox{0.95\linewidth}{\centering\fontsize{8pt}{10pt}\selectfont\textbf{Fuente:} elaboración propia.}
\end{minipage}
\vskip0.2cm



Con la función \verb"InterpolatingPolynomial" se consiguen dos animaciones: una con el gráfico y la otra con 
las expresiones que lo definen. \index{InterpolatingPolynomial}

\begin{tcolorbox}[breakable,enhanced,title=\bf Instrucción Mathematica] 
\begin{Verbatim}[formatcom=\letraverbatim]
Manipulate[Expand[InterpolatingPolynomial[{{1,1},
{7/5,k*6/5},{3,2}},x]],{k,12/10,2,1/10}]
Manipulate[Plot[Expand[InterpolatingPolynomial[{{1,1},
{7/5,k*6/5},{3,2}},x]],{x,1,3}],{k,-2,2,0.1}]
\end{Verbatim}
\end{tcolorbox}


\vspace{0pt}
\textcolor{black}{\scriptsize\ding{110}}\quad Utilizando las funciones $f(t)=\cos(\frac{\pi}{2}t)$ y $g(t)=\sin(\frac{\pi}{2}t)$, se define
$$\ve{r}(t)=\cos\left( \frac{\pi t}{2}\right)\left(1,1\right) +\sin \left( \frac{\pi t}{2}\right) \left( 3,2\right),$$
que permite unir los puntos $A(1,1)$ y $B(3,2).$
Esta expresión se consigue, al simplificar la expresión
$$\cos\left( \frac{\pi t}{2}\right)\left(1,1\right) +\sin \left( \frac{\pi t}{2}\right) \left( 3,2\right)$$
para $\,\,0\leq t\leq 1$. 
Por tanto, la parametrización escogida es $\ve{r}(t).$


\vskip0.2cm
\noindent
\begin{minipage}{\textwidth}
	\centering
	\captionof{figure}{${\bf r}(t)=\cos\left( \frac{\pi t}{2}\right)A +\sin \left( \frac{\pi t}{2}\right)B$}
	\begin{minipage}{0.28\linewidth}
		\centering
		\includegraphics[width=\linewidth]{Fig/31.pdf}
	\end{minipage}
	
	\vspace{0.3em}
	\parbox{0.95\linewidth}{\centering\fontsize{8pt}{10pt}\selectfont\textbf{Fuente:} elaboración propia.}
\end{minipage}
\vskip0.2cm


\section{Superficies} \index{Parametrización!de superficies} \label{superficies}
En esta sección construiremos algunas parametrizaciones de superficies a partir 
de los elementos usados en el tratamiento de curvas.
\subsection*{Parametrización de regiones planas}

Con las ideas elementales utilizadas para construir curvas, se pueden construir parametrizaciones de regiones en el plano.


\vskip0.2cm
\noindent
\begin{minipage}{\textwidth}
	\centering
	\captionof{figure}{${\bf u}(r,\theta) =\la r \cos \theta, r \sin \theta \ra$}
	\label{fig:32}
	\begin{minipage}{0.30\linewidth}
		\centering
		\includegraphics[width=\linewidth]{Fig/32.pdf}
	\end{minipage}
	
	\vspace{0.3em}
	\parbox{0.95\linewidth}{\centering\fontsize{8pt}{10pt}\selectfont\textbf{Fuente:} elaboración propia.}
\end{minipage}
\vskip0.2cm


Con el propósito de \emph{rellenar} una circunferencia de radio $a$, se puede variar el radio y el arco en las expresiones
$x= r \cos \theta$ y $y= r \sin \theta.$ Este círculo puede parametrizarse mediante la función $\ve{u}(r,\theta)= \la r \cos \theta, r \sin \theta \ra$,
en donde $r \in [0,a],$ $\theta \in [0,2 \pi].$

\vskip0.2cm

La expresión $4x^2+9y^2=36$ representa una elipse, esta curva se puede parametrizar mediante 
las ecuaciones $x=3\cos \theta,$ $y=2\sin \theta$.
La región interior se puede describir
mediante la función
${\bf u}(r,\theta) = \la 3 r \cos \theta,2r \sin\theta \ra,$
con $r\in [0,1],$ $\theta \in [0,2\pi].$


\vskip0.2cm
\noindent
\begin{minipage}{\textwidth}
	\centering
	\captionof{figure}{${\bf u}(r,\theta) =\la 3r \cos \theta, 2r \sin \theta \ra$}
	\label{fig:33}
	\begin{minipage}{0.26\linewidth}
		\centering
		\includegraphics[width=\linewidth]{Fig/33.pdf}
	\end{minipage}
	
	\vspace{0.3em}
	\parbox{0.95\linewidth}{\centering\fontsize{8pt}{10pt}\selectfont\textbf{Fuente:} elaboración propia.}
\end{minipage}
\vskip0.2cm


\vspace{0pt}
El triángulo con vértices $A$, $B$ y $C$, se parametriza a través de segmentos rectilíneos.
Se construye una combinación convexa entre un punto del segmento $\overline{AB}$ con el punto $C$.
Una función que describe esta región, con $u \in [0,1],$ $v \in [0,1],$ se obtiene simplificando la expresión


\vskip0.2cm
\noindent
\begin{minipage}{\textwidth}
	\centering
	\captionof{figure}{Triángul parametrizado $(1-v)C+v\left((1-u)A+uB \right).$}
	\label{fig:161}
	\begin{minipage}{0.38\linewidth}
		\centering
		\includegraphics[width=\linewidth]{Fig/161.pdf}
	\end{minipage}
	
	\vspace{0.3em}
	\parbox{0.95\linewidth}{\centering\fontsize{8pt}{10pt}\selectfont\textbf{Fuente:} elaboración propia.}
\end{minipage}
\vskip0.2cm


$$(1-v)C+v\left((1-u)A+uB \right).$$


El cuadrilátero con vértices $A(-2,-1),$ $B(-1,4)$, $C(3,1),$ $D(4,3),$ 
se parametriza formando dos combinaciones convexas adecuadamente elegidas (evitar crear un \emph{corbatín}) 
con el mismo parámetro; con las expresiones obtenidas se forma otra combinación 
convexa con un parámetro distinto. 


\vskip0.2cm
\noindent
\begin{minipage}{\textwidth}
	\centering
	\captionof{figure}{Cuadrilátero parametrizado}
	\label{fig:Cuadrilátero_parametrizado}
	\begin{minipage}{0.32\linewidth}
		\centering
		\includegraphics[width=\linewidth]{Fig/35.pdf}
	\end{minipage}
	
	\vspace{0.3em}
	\parbox{0.95\linewidth}{\centering\fontsize{8pt}{10pt}\selectfont\textbf{Fuente:} elaboración propia.}
\end{minipage}
\vskip0.2cm

Combinando las expresiones $(1-t)A+tB$ y $(1-t)C+tD$,
se obtiene la función $\ve{r}(t,u)= (1-u)\left ((1-t)A+tB\right )+u \left( (1-t)C+tD\right),$
con $t\in [0,1]$, $u\in [0,1].$ En este caso, %La función obtenida es 
$\ve{r}(t,u)=\la t+5u-2, 5t+2u-3tu -1\ra.$

Con las siguientes instrucciones se consigue la figura \ref{fig:Cuadrilátero_parametrizado}:
\par
\vskip -12pt

\begin{tcolorbox}[breakable,enhanced,title=\bf Instrucción \verb"Mathematica"] 
\begin{Verbatim}[formatcom=\letraverbatim]
ParametricPlot[(1-u)*((1-t)*{-2,-1}+t*{-1,4})
+u*((1-t)*{3,1}+t*{4,3}),{t,0,1},{u,0,1}]
\end{Verbatim}
\end{tcolorbox}


\vspace{0pt}
La región $R$ acotada por las figuras de las funciones $y=f(x)$ y $y=g(x)$,
con $x\in [a,b];$ se parametriza formando una combinación convexa entre dos puntos
$P(x,f(x))$ y $Q(x,g(x)).$ Por tanto, 
una función que describe a $R$ es
$\ve{r}(x,u)= (1-u)(x,f(x))+u(x,g(x)),$
con $u \in [0,1]$, $x \in [a,b].$


\vskip0.2cm
\noindent
\begin{minipage}{\textwidth}
	\centering
	\captionof{figure}{Región parametrizada}
	\label{fig:Región_parametrizada}
	\begin{minipage}{0.32\linewidth}
		\centering
		\includegraphics[width=\linewidth]{Fig/36.pdf}
	\end{minipage}
	
	\vspace{0.3em}
	\parbox{0.95\linewidth}{\centering\fontsize{8pt}{10pt}\selectfont\textbf{Fuente:} elaboración propia.}
\end{minipage}
\vskip0.2cm


De manera análoga al caso anterior, si la región está acotada por las figuras de las expresiones 
$x=u(y)$ y $x=v(y)$, con $y\in [c,d];$ con lo cual, para
$u \in [0,1]$, $y \in [c,d],$ 
una función que describe la dicha región es %\\[0.1cm]
\[
\ve{r}(y,u)= (1-u)\la y,u(y)\ra + u\la y,v(y)\ra
\]

\vskip0.2cm
\noindent
\begin{minipage}{\textwidth}
	\centering
	\captionof{figure}{Región parametrizada}
	\label{fig:Región_parametrizada_3}
	\begin{minipage}{0.32\linewidth}
		\centering
		\includegraphics[width=\linewidth]{Fig/37.pdf}
	\end{minipage}
	
	\vspace{0.3em}
	\parbox{0.95\linewidth}{\centering\fontsize{8pt}{10pt}\selectfont\textbf{Fuente:} elaboración propia.}
\end{minipage}
\vskip0.2cm


\vspace{0pt}
La región $R,$ acotada por las figuras de $y=4-x^2$ y $y=x^3-4x$ para $x\in[-1,2]$
se parametriza simplificando la expresión 
$(1-u)\la x, 4-x^2 \ra +u\la x, x^3-4x \ra $, es decir 
$${\bf r}(x,u)=\la x, (x^2-4)(ux+u-1) \ra,$$ con $x\in[-1,2],$ $u\in[0,1],$
representa a la mencionada región.


\vskip0.2cm
\noindent
\begin{minipage}{\textwidth}
	\centering
	\captionof{figure}{$y=4-x^2$ y $y=x^3-4x$}
	\label{fig:38}
	\begin{minipage}{0.28\linewidth}
		\centering
		\includegraphics[width=\linewidth]{Fig/38.pdf}
	\end{minipage}
	
	\vspace{0.3em}
	\parbox{0.95\linewidth}{\centering\fontsize{8pt}{10pt}\selectfont\textbf{Fuente:} elaboración propia.}
\end{minipage}
\vskip0.2cm


\vspace{0pt}
Con las ecuaciones $y=x+1$ y $y=x^2-1,$ para $x\in[-1,2],$
se forma la combinación convexa $(1-u)\la x, 1+x \ra +u\la x, x^2-1 \ra;$
simplificando se define la función ${\bf r}(x,u)=\la x, (x+1) (1-2u+ux) \ra,$ con $x\in[-1,2],$ $u\in[0,1],$
que determina la región mostrada.
	
\vskip0.2cm
\noindent
\begin{minipage}{\textwidth}
	\centering
	\captionof{figure}{$y=x+1$ y $y=x^2-1$}
	\label{fig:39}
	\begin{minipage}{0.33\linewidth}
		\centering
		\includegraphics[width=\linewidth]{Fig/39.png}
	\end{minipage}
	
	\vspace{0.3em}
	\parbox{0.95\linewidth}{\centering\fontsize{8pt}{10pt}\selectfont\textbf{Fuente:} elaboración propia.}
\end{minipage}
\vskip0.2cm	
	

\vspace{0pt}
La región acotada por las graficas de las expresiones $y=x^4-x^2,$ $y=4-x^2$ y $x=1$, 
se parametriza por medio de la expresión
$(1-u)(x,x^4-x^2)+u(x,4-x^2),$
que al simplificarse, con $x \in [-\sqrt{2},1]$, $u \in [0,1],$ permite definir la función \\
$\ve{r}(x,u)=\la x, u(4-x^4)+x^2(x^2-1) \ra.$


\vskip0.2cm
\noindent
\begin{minipage}{\textwidth}
	\centering
	\captionof{figure}{$\la x, u(4-x^4)+x^2(x^2-1) \ra$}
	\label{fig:region_gris}
	\begin{minipage}{0.35\linewidth}
		\centering
		\includegraphics[width=\linewidth]{Fig/40.pdf}
	\end{minipage}
	
	\vspace{0.3em}
	\parbox{0.95\linewidth}{\centering\fontsize{8pt}{10pt}\selectfont\textbf{Fuente:} elaboración propia.}
\end{minipage}
\vskip0.2cm	


Con las siguientes instrucciones se obtiene la figura \ref{fig:region_gris} en \verb"Mathematica"



\begin{tcolorbox}[breakable,enhanced,title=\bf Instrucción \verb"Mathematica"]  
\begin{Verbatim}[formatcom=\letraverbatim]
B=ParametricPlot[{x,u*(4-x^4)+x^2*(x^2-1)},
{x,-Sqrt[2],1},{u,0,1}]; 
A=ContourPlot[{y==x^4-x^2,y==4-x^2},{x,-2,2},{y,-1,4}];
Show[A,B]
\end{Verbatim}
\end{tcolorbox}


\vspace{0pt}
Las parábolas $x=y^2-5y,$ $x=3y-y^2$ se cortan 
en $(0,0)$ y $(-4,4).$ La región acotada por estas curvas se parametriza 
simplificando la expresión
$(1-v)\la u^2-5u,u\ra +v\la 3u-u^2,u\ra,$
con $u \in [0,4]$, $v \in [0,1]$. Esta región se describe por medio de la función
$\ve{r}(u,v) = \la u (u+8v-2uv-5),u \ra.$


\vskip0.2cm
\noindent
\begin{minipage}{\textwidth}
	\centering
	\captionof{figure}{$ \la (8v-5)u+(1-2v)u^2,u \ra$}
	\label{fig:41}
	\begin{minipage}{0.27\linewidth}
		\centering
		\includegraphics[width=\linewidth]{Fig/41.pdf}
	\end{minipage}
	
	\vspace{0.3em}
	\parbox{0.95\linewidth}{\centering\fontsize{8pt}{10pt}\selectfont\textbf{Fuente:} elaboración propia.}
\end{minipage}
\vskip0.2cm	


\begin{tcolorbox}[breakable,enhanced,title=\bf Instrucción \verb"Mathematica"] 
\begin{Verbatim}[formatcom=\letraverbatim]
A=ContourPlot[{x==y^2-5*y,x==3*y-y^2},{x,-7,4},
{y,-1,5}];
B=ParametricPlot[(1-v){u^2-5*u,u}+v*{3*u-u^2,u},
{v,0,1},{u,0,4}];
Show[A,B]
\end{Verbatim}
\end{tcolorbox}

\subsection*{Superficies parametrizadas} \index{Superficies parametrizadas}
Recordemos un concepto que permite hacer estas construcciones en el espacio.

\begin{tcolorbox}[breakable,enhanced,title=\bf Definición]
Una parametrización de una superficie $\sigma $,
es una aplicación $$\ve{r}:[a,b]\times \lbrack c,d]\rightarrow \mathbb{R}^{3},$$
tal que al variar los parámetros $u\in \lbrack a,b]$ y $v\in \lbrack c,d]$
su imagen $\ve{r}(u,v)$ va describiendo los puntos de la superficie $\sigma $.
\end{tcolorbox}

\textcolor{black}{\scriptsize\ding{110}}\quad Utilizando coordenadas esféricas, la superficie de la esfera unitaria se puede parametrizar 
por medio de la función
{\small
\[{\bf r}(\theta, \phi)=\big \langle \cos \theta \sin \phi, \sin \theta \sin \phi, \cos \phi \big\rangle,\]
}
con $\theta \in [0,2\pi]$ y $ \phi \in [0,\pi]$.

\vskip0.2cm

El elipsoide con ecuación $\frac{x^2}{16}+\frac{y^2}{9}+\frac{z^2}{4}=1,$ 
se parametriza con una sencilla variación de la función anterior
${\bf r}(\theta, \phi)=\la 4\cos \theta \sin \phi,3\sin \theta \sin \phi, 2\cos \phi \ra.$ 

\vskip0.2cm
\noindent
\begin{minipage}{\textwidth}
	\centering
	\captionof{figure}{\small La esfera unitaria y el elipsoide parametrizado por ${\bf r}(\theta, \phi)$}
	\begin{minipage}{0.20\linewidth}
		\centering
		\includegraphics[width=\linewidth]{Img/T_53a.pdf}
	\end{minipage}
	\begin{minipage}{0.30\linewidth}
		\centering
		\includegraphics[width=\linewidth]{Img/T_53b.pdf}
	\end{minipage}
	\begin{minipage}{0.30\linewidth}
		\centering
		\includegraphics[width=\linewidth]{Img/T_53c.pdf}
	\end{minipage}
	
	\vspace{0.3em}
	\parbox{0.95\linewidth}{\centering\fontsize{8pt}{10pt}\selectfont\textbf{Fuente:} elaboración propia.}
\end{minipage}
\vskip0.2cm


Esta figura se obtiene al ejecutar las siguientes instrucciones:

\begin{tcolorbox}[breakable,enhanced,title=\bf Instrucción \verb"Mathematica"]  
\begin{Verbatim}[formatcom=\letraverbatim]
ParametricPlot3D[{Cos[\[Theta]]*Sin[\[Phi]],
Sin[\[Theta]]*Sin[\[Phi]],Cos[\[Phi]]},{\[Theta],0,
2*Pi},{\[Phi],0,Pi},MeshStyle->Gray,PlotStyle->Orange]

ParametricPlot3D[{4Cos[\[Theta]]*Sin[\[Phi]],3*
Sin[\[Theta]]*Sin[\[Phi]],2*Cos[\[Phi]]},{\[Theta],0,
2*Pi},{\[Phi],0,Pi},MeshStyle->Gray,PlotStyle->Orange]
\end{Verbatim}
\end{tcolorbox}


%\vspace{-5mm}
\textcolor{black}{\scriptsize\ding{110}}\quad La superficie $\sigma$ sobre el paraboloide 
$2z=9-x^2-y^2$ y dentro del cilindro $x^2+ \left(y -\frac{1}{3} \right)^2=4,$ se parametriza al asignar $x=u \cos v,$ 
$y= \frac{1}{3}+u\sin v$, con $u \in [0,2]$, $v \in [0,2\pi].$ 
Del paraboloide se obtiene $z = \frac{1}{18} (80-9u^2-6u\sin(t)).$


\vskip0.2cm
\noindent
\begin{minipage}{\textwidth}
	\centering
	\captionof{figure}{Superficie $\sigma$}
	\label{Superficie_sigma}
	\begin{minipage}{0.38\linewidth}
		\centering
		\includegraphics[width=\linewidth]{Fig/42.pdf}
	\end{minipage}
	
	\vspace{0.3em}
	\parbox{0.95\linewidth}{\centering\fontsize{8pt}{10pt}\selectfont\textbf{Fuente:} elaboración propia.}
\end{minipage}
\vskip0.2cm


Entonces, para $u \in [0,2]$, $v \in [0,2\pi],$ la superficie $\sigma$ se parametriza por
\[{\bf r}(u,v)= \langle u\cos v, \frac{1}{3} +u\sin v, \frac{1}{18}(80-9u^2-6u \sin(v)) \rangle.\]
Si $u=2$ la expresión obtenida corresponde a la parametrización de la curva de intersección del paraboloide
y el cilindro, es decir ${\bf r}(v)= \langle 2\cos v, \frac{1}{3} +2\sin v,\frac{1}{9}(22-6\sin(v))\rangle,$ $v \in [0,2\pi].$


\vspace{0pt}
\textcolor{black}{\scriptsize\ding{110}}\quad La superficie del cilindro con ecuación $x^2+y^2=2x$, para $z \in [-1,2],$ se 
parametriza en la forma 
$${\bf r}(u,t)= \la 1+\cos(t),\sin(t), u \ra,$$ con $t \in [0,2\pi],$ $u\in [-1,2].$ 
En el plano $xy$ el cilindro corresponde a la circunferencia $(x-1)^2+y^2=1.$


\vskip0.2cm
\noindent
\begin{minipage}{\textwidth}
	\centering
	\captionof{figure}{${\bf r}(u,t)= \la 1+\cos t,\sin t, u \ra$}
	\label{fig:43}
	\begin{minipage}{0.32\linewidth}
		\centering
		\includegraphics[width=\linewidth]{Fig/43.pdf}
	\end{minipage}
	
	\vspace{0.3em}
	\parbox{0.95\linewidth}{\centering\fontsize{8pt}{10pt}\selectfont\textbf{Fuente:} elaboración propia.}
\end{minipage}
\vskip0.2cm


\textcolor{black}{\scriptsize\ding{110}}\quad Las expresiones $2z=x^{2}-y^{2}$ y $2x^{2}+y^{2}+2y=8$ representan un paraboloide hiperbólico y un cilindro elíptico, respectivamente. La ecuación
del cilindro representa una elipse en el plano $xy$, el interior de esta región 
se parametriza como $x=\frac{3}{\sqrt{2}}u\cos (t),$ $y=-1+3u\sin (t),$ con la ecuación del paraboloide hiperbólico 
se halla $z$, esto es 
$z=\frac{1}{2}\left( x^{2}-y^{2}\right) =\frac{27}{4}u^{2}\cos^{2}(t) +3u\sin (t) -\frac{9}{2}u^{2}-\frac{1}{2}.$ 
Entonces, la función
{\small
$${\bf r}\left( t,u\right) =\left\langle \frac{3}{\sqrt{2}}u\cos (t),3u\sin (t)-1 ,\frac{27}{4}u^{2}\cos ^{2}(t)
+3u\sin (t) -\frac{9u^{2}+1}{2} \right\rangle,$$
}

con $t\in \lbrack 0,2\pi ],$ $u\in \lbrack 0,1],$  parametriza la región sobre el paraboloide 
hiperbólico acotada por el cilindro.
\vskip0.2cm
\noindent
\begin{minipage}{\textwidth}
	\centering
	\captionof{figure}{\small Superficie sobre el paraboloide hiperbólico.}
	\begin{minipage}{0.23\linewidth}
		\centering
		\includegraphics[width=\linewidth]{Img/73a.pdf}
	\end{minipage}
	\hspace{1em}
	\begin{minipage}{0.23\linewidth}
		\centering
		\includegraphics[width=\linewidth]{Img/73b.pdf}
	\end{minipage}
	
	\vspace{0.3em}
	\parbox{0.95\linewidth}{\centering\fontsize{8pt}{10pt}\selectfont\textbf{Nota:} elaboración propia.}
\end{minipage}
\vskip0.2cm

La anterior figura se consigue con las siguientes instrucciones.

\begin{tcolorbox}[breakable,enhanced,title=\bf Instrucción \verb"Mathematica"]  
\begin{Verbatim}[formatcom=\letraverbatim]
a=ContourPlot3D[{2*z==x^2-y^2,2*x^2+y^2+2*y==8},
{x,-5,5},{y,-5,5},{z,-9,5},Mesh->False,ContourStyle
->{Directive[Orange,Opacity[0.6]],Directive[Cyan,
Opacity[0.2]]},PlotPoints->80,Ticks->{{-3,0,3},
{-3,0,3},{-8,-4,-2,2,4}}];
b=ParametricPlot3D[{3*u*Cos[t]/Sqrt[2],-1+3*u*Sin[t],
27*u^2*Cos[t]^2/4+3*u*Sin[t]-9*u^2/2-1/2},{t,0,2*Pi},
{u,0,1},PlotStyle->{Cyan},PlotPoints->80];
Show[a,b]
\end{Verbatim}
\end{tcolorbox}

\textcolor{black}{\scriptsize\ding{110}}\quad Las ecuaciones $x^{2}+\left( y-3\right) ^{2}=4$ y $4z=36-\left(x^{2}+y^{2}\right) $ 
representan un cilindro y un paraboloide invertido con vértice en el punto $(0,0,9)$.
Con el fin de parametrizar la superficie $\sigma$ sobre el paraboloide exterior al cilindro, con $z\geq 0,$
analizaremos la base y luego proyectaremos esta región sobre el paraboloide.
La figura \ref{Superficie_paraboloide_exterior_cilindro} del sólido en mención se obtiene con ayuda de las siguientes instrucciones:

\begin{tcolorbox}[breakable,enhanced,title=\bf Instrucción \verb"Mathematica"] 
\begin{Verbatim}[formatcom=\letraverbatim]
ContourPlot3D[{z==9-(x^2+y^2)/4,x^2+(y-3)^2==4},
{x,-6,6},{y,-6,6},{z,0,9},Mesh->False,ContourStyle
->{Orange,Cyan}]
ParametricPlot3D[{2*(2*u+1)*Sin[t],2*(2*u+1)Cos[t]
-3*u+3,(u-1)*(12*(1+2u)*Cos[t]-23-25u)/4},{u,0,1},
{t,0,2*Pi}]
\end{Verbatim}
\end{tcolorbox}


\noindent
\begin{minipage}{\textwidth}
	\centering
	\captionof{figure}{Superficie sobre el paraboloide exterior al cilindro}
	\label{Superficie_paraboloide_exterior_cilindro}
	\begin{minipage}{0.26\linewidth}
		\centering
		\includegraphics[width=\linewidth]{Img/86a.pdf}
	\end{minipage}
	\hspace{1em}
	\begin{minipage}{0.33\linewidth}
		\centering
		\includegraphics[width=\linewidth]{Fig/44.pdf}
	\end{minipage}
	
	\vspace{0.3em}
	\parbox{0.95\linewidth}{\centering\fontsize{8pt}{10pt}\selectfont\textbf{Fuente:} elaboración propia.}
\end{minipage}
\vskip0.2cm





La base del cilindro se puede parametrizar como
\begin{equation} \label{sup5a}
x=2\sin (t), \,\, y=3+2\cos (t), \, t\in \lbrack 0,2\pi ].
\end{equation}

Mientras que la intersección del paraboloide y el plano $z=0$ se parametriza como $x=6\sin (t) $, $y=6\cos (t) $, con $t\in \lbrack 0,2\pi ].$


Con el fin de parametrizar la proyección de $\sigma$ sobre el plano $xy$, se puede formar 
una combinación convexa entre dos puntos de la forma 
$P = \left( 2\sin(t),\, 3 + 2\cos(t),\, 0 \right)$ y 
$Q = \left( 6\sin(t),\, 6\cos(t),\, 0 \right)$. 
La expresión que permite describir la región acotada por las dos circunferencias está dada por 
\[
(1 - u)\,P + u\,Q.
\]

\vskip0.2cm
\noindent
\begin{minipage}{\textwidth}
	\centering
	\captionof{figure}{Proyección de $\sigma$}
	\label{fig:45}
	\begin{minipage}{0.33\linewidth}
		\centering
		\includegraphics[width=\linewidth]{Fig/45.pdf}
	\end{minipage}
	
	\vspace{0.3em}
	\parbox{0.95\linewidth}{\centering\fontsize{8pt}{10pt}\selectfont\textbf{Fuente:} elaboración propia.}
\end{minipage}
\vskip0.2cm

Por tanto, la función que parametriza la región sobre el plano $xy$ acotada por las dos 
circunferencias presentadas antes es
\begin{equation}
\ve{v}(u,t)=\left( 2\left( 2u+1\right) \sin t,2\left( 2u+1\right) \cos t-3u+3,0\right) \label{sup5}
\end{equation}
La coordenada $z$ sobre la superficie del paraboloide es
\begin{eqnarray*}
z &=&9-\frac{1}{4}\left( \left( 2\left( 2u+1\right) \sin t\right) ^{2}+\left(
2\left( 2u+1\right) \cos t-3u+3\right) ^{2}\right) \\
&=& \frac{1}{4}\left( 1-u\right) \left( 25u-12(2u+1) \cos t+23\right)
\end{eqnarray*}
Por lo anterior, la superficie $\sigma$ se parametriza mediante la función 
\begin{multline*}
{\bf r}(t,u) =\big\langle 2(2u+1) \sin t,2(2u+1)\cos t-3u+3, \\
\frac{1}{4}( 1-u) ( 25u-12(2u+1)\cos t+23)  \big\rangle,
\end{multline*}
donde $t \in [0,2\pi],$ $u \in [0,1].$
Sobre la intersección de las superficies se satisface que 
$z=9-\frac{1}{4}\left( \left( 2\sin(t)\right)^{2}+\left(3+2\cos(t)\right)^{2}\right) = \frac{23}{4}-3\cos(t);$
dicha curva se parametriza como 
$$\ve{v}(t)= \langle 2\sin (t), 3+2\cos (t), \frac{23}{4}-3\cos(t) \rangle, \, t \in [0,2\pi].$$
\vskip0.3cm
\textcolor{black}{\scriptsize\ding{110}}\quad Con algunos de los elementos anteriores construiremos el siguiente ejemplo.
Utilizando la ecuación (\ref{sup5}), se sigue que
la coordenada $z$ de la curva de intersección del cilindro $x^2+(y-3)^2=4$ y el paraboloide hiperbólico 
$5z=x^2-y^2$ está dado por
{\small
\[
z=\frac{x^{2}-y^{2}}{5}=\frac{1}{5}\left((2\sin(t))^{2}-(3+2\cos(t))^{2}\right)=-1-\frac{12}{5}\cos t-\frac{8}{5}\cos ^{2}t.
\]
}

La región determinada por la ecuación (\ref{sup5}), se proyecta sobre el paraboloide hiperbólico con ecuación
$5z=x^2-y^2;$ la coordenada $z$ está dada por
{\small
\begin{eqnarray*}
z &=& \frac{1}{5}\left( x^{2}-y^{2}\right) =\frac{1}{5}\left( \left( 2\left(1+2u\right) \sin t\right) ^{2}-\left( 3-3u+\left( 2+4u\right) \cos t\right)^{2}\right) \\
&=& \frac{1}{5}\left( u+5\right) \left( 7u-1\right) -\frac{8}{5}\left( 1+2u\right) ^{2}\cos ^{2}t+\frac{12}{5}\left( 1+2u\right) \left(u-1\right) \cos t
\end{eqnarray*}
}

Entonces, la función
{\small
\begin{multline*}
{\bf r}(t,u) = \big \langle 2\left( 1+2u\right) \sin t,3-3u+\left( 2+4u\right) \cos t, \frac{1}{5} (u+5) (7u-1)\\
-\frac{8}{5}\left( 1+2u\right) ^{2}\cos ^{2}t+\frac{12}{5}\left( 1+2u\right) \left( u-1\right) \cos t\big \rangle, 
\end{multline*}
}
$u\in \lbrack 0,1],$ $t\in \lbrack 0,2\pi ],$ parametriza la superficie sobre
la silla de montar con ecuación  $5z=x^{2}-y^{2},$ exterior al cilindro con ecuación 
$x^{2}+(y-3)^{2}=4,$ e interior al paraboloide con ecuación 
$z=9-\left( x^{2}-y^{2}\right) /4.$
La figura \ref{Superficie_paraboloide_hiper} ilustra dicha región.

\noindent
\begin{minipage}{\textwidth}
	\centering
	\captionof{figure}{Superficie sobre el paraboloide hiperbólico $5z=x^2-y^2$}
	\label{Superficie_paraboloide_hiper}
	\begin{minipage}{0.47\linewidth}
		\centering
		\includegraphics[width=\linewidth]{Fig/46.pdf}
	\end{minipage}
	
	\vspace{0.3em}
	\parbox{0.95\linewidth}{\centering\fontsize{8pt}{10pt}\selectfont\textbf{Fuente:} elaboración propia.}
\end{minipage}
\vskip0.2cm

La figura de la superficie se obtiene ejecutando las siguientes instrucciones. Este es un ejemplo de parametrización en \verb"Mathematica".

\begin{tcolorbox}[breakable, enhanced, title={Código en \verb"Mathematica"}]
\begin{Verbatim}[formatcom=\letraverbatim]
{x,y}=(1-u)*{2*Sin[t],3+2*Cos[t]}+u*{6*Sin[t],6*Cos[t]};
r={x,y,(x^2-y^2)/5};
ParametricPlot3D[r,{u,0,1},{t,0,2Pi},ViewPoint->{-3,6,2}]
\end{Verbatim}
\end{tcolorbox}

Con los elementos anteriores es posible parametrizar otras superficies.
\begin{tcolorbox}[breakable,enhanced,title=\bf Una manera de parametrizar una superficie]
Para dos funciones vectoriales $\ve{\alpha}$ y $\ve{\beta}$ en $\mathbb{R}^n$, si $f$ y $g$ son
funciones reales contínuas en el intervalo $[a,b]$, que satisfacen que 
$f(a)=1,$ $g(a)=0$ y $f(b)=0,$ $g(b)=1,$ entonces la función vectorial
$$\ve{r}(u,t)=f(u)\ve{\alpha}(t)+g(u)\ve{\beta}$$
es una parametrización de la superficie $\sigma$ acotada por las curvas determinadas por $\ve{\alpha}$ y $\ve{\beta}$.
\end{tcolorbox}

Veamos algunos ejemplos de la expresión anterior.

\textcolor{black}{\scriptsize\ding{110}}\quad La cardioide $r=1-\cos(t),$ en el plano $z=2,$ se define por medio de la función vectorial 
${\bf a}(t)= \la (1-\cos(t))\cos(t),(1-\cos(t))\cos(t),2  \ra,$  con $t \in [0,2\pi],$
mientras que la función 
${\bf b}(t)= \la 3\cos(t), 3\sin(t), 0\ra,$ con $t \in [0,2\pi],$ 
determina una circunferencia en el plano $xy.$ 
Una combinación convexa de estas dos funciones vectoriales determina, mediante segmentos rectilíneos,
la superficie que se muestra a continuación. \index{Cardioide} 
\index{Combinación convexa}

\vskip0.2cm
\noindent
\begin{minipage}{\textwidth}
	\centering
	\captionof{figure}{Superficie que conecta una cardioide y una circunferencia}
	\label{Superficie_cardioide_circunferencia}
	\begin{minipage}{0.21\linewidth}
		\centering
		\includegraphics[width=\linewidth]{Fig/47.pdf}
	\end{minipage}
	\hspace{1em}
	\begin{minipage}{0.23\linewidth}
		\centering
		\includegraphics[width=\linewidth]{Img/42a.pdf}
	\end{minipage}
	\hspace{1em}
	\begin{minipage}{0.25\linewidth}
		\centering
		\includegraphics[width=\linewidth]{Img/42b.pdf}
	\end{minipage}
	
	\vspace{0.3em}
	\parbox{0.95\linewidth}{\centering\fontsize{8pt}{10pt}\selectfont\textbf{Fuente:} elaboración propia.}
\end{minipage}
%\vskip0.2cm
%\clearpage

Simplificando la expresión $(1-u){\bf a}(t)+u{\bf b}(t),$
se obtiene la función que parametriza esta superficie
{\small
\[{\bf r}(u,t)= \langle \cos(t)((u-1)\cos(t)+1+2u), \sin(t)((u-1)\cos(t)+1+2u),2-2u\rangle,\]
}
con $t\in[0,2\pi],\, u\in[0,1],$ la figura anterior se obtiene al ejecutar las siguientes instrucciones:


\begin{tcolorbox}[breakable,enhanced,title=\bf Instrucción \verb"Mathematica"] 
\begin{Verbatim}[formatcom=\letraverbatim]
q={(1-Cos[t])*Cos[t],(1-Cos[t])*Sin[t],2};
r=(1-u)*q+u*{3*Cos[t],3*Sin[t],0};
a=ParametricPlot3D[{q,{3*Cos[t],3*Sin[t],0}},{t,0,
2*Pi},PlotStyle->{Directive[Black,Thickness->0.01],
PlotPoints->90,Directive[Black,Thickness->0.01]}];
b=ParametricPlot3D[r,{t,0,2*Pi},{u,0,1},
PlotStyle->Orange,MeshStyle->Gray];
Show[a,b]
\end{Verbatim}
\end{tcolorbox}

\textcolor{black}{\scriptsize\ding{110}}\quad Con las funciones vectoriales
$\overrightarrow{r}_1(t)=\la 3\cos(t),\sin (t),4+\sin ( t) \ra $ y
$\overrightarrow{r}_2(t)=\la 2\cos(t),3\sin (t),-1+\cos ( t) \ra $ con $ t\in \lbrack 0,2\pi],$
que representan dos elipses en el espacio, se forma una combinación convexa de ellas con parámetro 
$u\in \lbrack 0,1].$
\\
Simplificando la expresión $(1-u)\overrightarrow{r}_1(t)+u\overrightarrow{r}_2(t)$
se obtiene la función:
$${\bf r}(u,t)=\la (3-u)\cos(t),(1+2u)\sin(t),4+(1-u)\sin(t)-5u+u\cos(t)\ra,$$
para $u\in[0,1],\,\, t\in[0,2\pi],$ la cual representa la superficie que se obtiene de unir las dos elipses mediante segmentos rectilíneos. 



\vskip0.2cm
\noindent
\begin{minipage}{\textwidth}
	\centering
	\captionof{figure}{\small Superficie que une dos elipses mediante segmentos}
	\label{conos_67}
	\begin{minipage}{0.17\linewidth}
		\centering
		\includegraphics[width=\linewidth]{Img/33a.pdf}
	\end{minipage}
	\begin{minipage}{0.17\linewidth}
		\centering
		\includegraphics[width=\linewidth]{Img/33b.pdf}
	\end{minipage}
	\begin{minipage}{0.17\linewidth}
		\centering
		\includegraphics[width=\linewidth]{Img/33c.pdf}
	\end{minipage}
	
	\vspace{0.3em}
	\parbox{0.95\linewidth}{\centering\fontsize{8pt}{10pt}\selectfont\textbf{Fuente:} elaboración propia.}
\end{minipage}
\vskip0.2cm


La figura \ref{conos_67} se consigue ejecutando las siguientes instrucciones:

\begin{tcolorbox}[breakable,enhanced,title=\bf Instrucción \verb"Mathematica"]  
\begin{Verbatim}[formatcom=\letraverbatim]
ParametricPlot3D[(1-u)*{3*Cos[t],Sin[t],4+Sin[t]}+
u*{2*Cos[t],3*Sin[t],-1+Cos[t]},{t,0,2*Pi},{u,0,1},
MeshStyle->Blue,PlotStyle->{Opacity[0.9]}]
\end{Verbatim}
\end{tcolorbox}

También se pueden generar dos figuras por separado llamadas \verb"a" y \verb"b", 
luego, con ayuda de la función \verb"Show", se muestran en una sola.

\begin{tcolorbox}[breakable,enhanced,title=\bf Instrucción \verb"Mathematica"]  
\begin{Verbatim}[formatcom=\letraverbatim]
a=ParametricPlot3D[{{3*Cos[t],Sin[t],4+Sin[t]},
{2*Cos[t],3*Sin[t],-1+Cos[t]}},{t,0,2*Pi},
PlotStyle->{Directive[RGBColor[0.2,0.1,1],
Thickness->0.015],Directive[Black,Thickness->0.015]}];
b=ParametricPlot3D[(1-u)*{3*Cos[t],Sin[t],4+Sin[t]}+
u*{2*Cos[t],3*Sin[t],-1+Cos[t],{t,0,0,2*Pi},{u,0,1},
MeshStyle->Gray,PlotStyle->{Orange,Opacity[0.9]}];
Show[a,b]
\end{Verbatim}
\end{tcolorbox}

%\fontsize{9.5}{12}\selectfont
\textcolor{black}{\scriptsize\ding{110}}\quad Las dos circunferencias en el espacio, representadas por las funciones vectoriales 
$\overrightarrow{\alpha}(t)=\langle 2\cos(t), 2\sin(t),0\rangle $ y $\overrightarrow{\beta}(t)=\langle \cos(t),\sin(t),2\rangle,$
con $t \in [0,2\pi].$ 
Con ellas se forma la siguiente combinación:
$$\frac{2\sqrt[4]{1-u}}{e^u+1}\, \overrightarrow{\alpha}(t) + \sqrt[4]{u}\,\overrightarrow{\beta}(t).$$
Con diferentes valores de $u= \frac{1}{4}, \,\frac{1}{2},\,2$, se logra la figura \ref{superficies_circunferencias68}


\begin{figure}[H]
	\begin{center}
		\captionof{figure}{\small Superficies que conectan dos circunferencias.}
		\label{superficies_circunferencias68}
		\includegraphics[width=2.75cm,height=1.5cm]{Img/8c.pdf}
		\includegraphics[width=2.75cm,height=1.5cm]{Img/8b.pdf}
		\includegraphics[width=2.75cm,height=1.5cm]{Img/8a.pdf}
	\end{center}
	\parbox{0.95\linewidth}{\centering\fontsize{8pt}{10pt}\selectfont\textbf{Fuente:} elaboración propia.}
\end{figure}


Con ayuda de la función \verb"Manipulate" se pueden generar varios elementos como los anteriores. 
Los comandos necesarios son los siguientes:

\begin{tcolorbox}[breakable,enhanced,title=\bf Instrucción \verb"Mathematica"] 
\begin{Verbatim}[formatcom=\letraverbatim]
Manipulate[ParametricPlot3D[2*(1-u)^k/(Exp[u]+1)
*{2*Cos[t],2*Sin[t],0}+u^k*{Cos[t],Sin[t],2},
{t,0,2Pi},{u,0,1}],{k,0.5,2,0.1}]
\end{Verbatim}
\end{tcolorbox}

\textcolor{black}{\scriptsize\ding{110}}\quad Con las funciones $f(t)=te^{t-1}$ y $g(t)=(1-t)e^{t}$, 
es sencillo ver que $f(0)=0=g(1)$ y $f(1)=1=g(0)$. Además, con las funciones vectoriales
\vspace{-2mm}
{\small
\[
\ve{\alpha}(t) = \left\langle 3\cos(t),2\sin(t),2-\frac{\sin(t)}{2}\right\rangle 
\, \text{y} \, 
\ve{\beta}(t) = \left\langle 2\cos(t),\sin(t),\frac{\sin (2t)}{2} \right\rangle
\]
}

se define la función  
${\bf r}(u,t)= f(u)\ve{\alpha}(t)+g(u)\ve{\beta}(t),$ $t\in [0,2\pi],$ $u\in [0,1],$ 
cuya figura es la siguiente:

\noindent
\begin{minipage}{\textwidth}
	\centering
	\captionof{figure}{Superficie que conecta las curvas $\alpha$ y $\beta$}
	\label{superficie_conecta_alfa_beta}
	\begin{minipage}{0.22\linewidth}
		\centering
		\includegraphics[width=\linewidth]{Fig/48.pdf}
	\end{minipage}
	\hspace{1em}
	\begin{minipage}{0.22\linewidth}
		\centering
		\includegraphics[width=\linewidth]{Fig/49.pdf}
	\end{minipage}
	
	\vspace{0.3em}
	\parbox{0.95\linewidth}{\centering\fontsize{8pt}{10pt}\selectfont\textbf{Fuente:} elaboración propia.}
\end{minipage}

\textcolor{black}{\scriptsize\ding{110}}\quad Las siguientes funciones vectoriales
{\small
$\ve{\alpha}(t)= \langle3\cos(t),2\sin(t),2-\frac{1}{2}\sin(t)\rangle$ y 
$\ve{\beta}(t)= \langle 2\cos(t),\sin(t),\frac{1}{2}\sin ( 2t) \rangle,$ $t\in [0,2\pi]$
} determinan dos curvas en el espacio. 
Con ayuda de las funciones $f(t)=1-t^2$ y $g(t)=t^3$, se puede formar la combinación entre ellas y
definir la función ${\bf r}(u,t)= f(u)\ve{\alpha}(t)+g(u)\ve{\beta}(t)$ con $t \in [0,2\pi],$ $u \in [0,1].$ 
Esta une los puntos de las curvas determinadas por $\ve{\alpha}(t)$
y $\ve{\beta}(t),$ a través de una superficie curva.

\vskip0.2cm
\noindent
\begin{minipage}{\textwidth}
	\centering
	\captionof{figure}{Superficie determinada por ${\bf r}(u,t)$}
	\label{Superficie_determinada}
	\begin{minipage}{0.23\linewidth}
		\centering
		\includegraphics[width=\linewidth]{Fig/50.pdf}
	\end{minipage}
	\hspace{1em}
	\begin{minipage}{0.30\linewidth}
		\centering
		\includegraphics[width=\linewidth]{Img/69a.pdf}
	\end{minipage}
	
	\vspace{0.3em}
	\parbox{0.95\linewidth}{\centering\fontsize{8pt}{10pt}\selectfont\textbf{Fuente:} elaboración propia.}
\end{minipage}
\vskip0.2cm

Esta figura se genera con ayuda de las siguientes instrucciones:

\begin{tcolorbox}[breakable,enhanced,title=\bf Instrucción \verb"Mathematica"] 
\begin{Verbatim}[formatcom=\letraverbatim]
a={3*Cos[t],2*Sin[t],2-Sin[t]/2};
b={2*Cos[t],Sin[t],Sin[2*t]/2};
p=ParametricPlot3D[(1-u^2)*a+u^3*b,{t,0,2*Pi},{u,0,1},
PlotStyle->Orange,MeshStyle->Gray,PlotPoints->60];
q=ParametricPlot3D[{a,b},{t,0,2*Pi},PlotPoints->
60,PlotStyle->{Directive[Black,Thickness->0.012],
Directive[Black,Thickness->0.012]}];
Show[p,q]
\end{Verbatim}
\end{tcolorbox}

\vskip0.2cm
\vspace{0pt}
\textcolor{black}{\scriptsize\ding{110}}\quad Las dos circunferencias unitarias ubicadas en el plano $z=2$ y $z=0$,
se unen por medio de curvas y generan la siguiente superficie.
La parametrizacion de esta región se consigue con la ayuda de las funciones 
$f(t)=\frac{2t}{t^2+1}$ y $g(t)=\frac{1-t}{t^2+1}$.


\vskip0.2cm
\noindent
\begin{minipage}{\textwidth}
	\centering
	\captionof{figure}{Circunferencias conectadas}
	\begin{minipage}{0.32\linewidth}
		\centering
		\includegraphics[width=\linewidth]{Fig/51.pdf}
	\end{minipage}
	
	\vspace{0.3em}
	\parbox{0.95\linewidth}{\centering\fontsize{8pt}{10pt}\selectfont\textbf{Fuente:} elaboración propia.}
\end{minipage}
\vskip0.2cm


Estas funciones se utilizan para formar la siguiente combinación.
{\small
\begin{equation*}
f(u)\ve{\alpha}(t)+g(u)\ve{\beta}(t) 
% \frac{2u}{u^2+1}\la \cos(t),\sin(t),2\ra + \frac{1-u}{u^2+1} \la\cos(t),\sin(t),0\ra\\
= \la  \frac{u+1}{u^2+1}\cos t,\frac{u+1}{u^2+1} \sin t,\frac{4u}{u^2+1}\ra,u \in[0,1], t\in[0,2\pi].
\end{equation*}
}

\vspace{0pt}
\textcolor{black}{\scriptsize\ding{110}}\quad Las funciones {\small $\ve{\alpha}(t) = \langle(2-2\cos(t)) \cos(t), (2-2\cos(t))\sin(t),2\rangle$} y {\small $\ve{\beta}(t) = \langle ( 1-\cos(t)) \cos(t),( 1-\cos(t)) \sin(t),0\rangle$} 
representan dos cardioides.

Se forma la combinación {\small $${\bf r}(u,t)= (1-u^{4})\ve{\alpha}(t) +u^{3}\ve{\beta}(t),$$}
con $u \in[0,1],$ $t \in[0,2\pi],$ que parametriza la superficie curva que las une y se muestra en la figura \ref{Dos_cardioides_conectadas}.


\vskip0.2cm
\noindent
\begin{minipage}{\textwidth}
	\centering
	\captionof{figure}{Dos cardioides conectadas}
	\label{Dos_cardioides_conectadas}
	\begin{minipage}{0.30\linewidth}
		\centering
		\includegraphics[width=\linewidth]{Fig/52.pdf}
	\end{minipage}
	
	\vspace{0.3em}
	\parbox{0.95\linewidth}{\centering\fontsize{8pt}{10pt}\selectfont\textbf{Fuente:} elaboración propia.}
\end{minipage}
\vskip0.2cm


\textcolor{black}{\scriptsize\ding{110}}\quad La figura \ref{Dos_cardioides_conectadas} de la función vectorial {\small ${\bf r}(t)= \langle 4 \cos(t)-\cos(4t),4\sin(t)-\sin(4t),0 \rangle,$} con
$t\in [0,2\pi],$ se presenta a continuación.

\vskip0.2cm
\vspace{0pt}
El interior se parametriza con
$x(u,t) = (4 \cos(t)-\cos(4t))u,$
$y(u,t) = (4\sin(t)-\sin(4t))u,$
para $t\in [0,2\pi],$ $u\in [0,1].$
Utilizando las funciones $f(u)= \sqrt[4]{(1-u)^3}$  y $g(u)=\sqrt[4]{u^3}$,
se forma la combinación con el punto $(1,0,3),$
$${\bf r}_2(u,t)=f(u) {\bf r}(t)+ g(u)(1,0,3).$$


\vskip0.2cm
\noindent
\begin{minipage}{\textwidth}
	\centering
	\captionof{figure}{${\bf r}_1(u,t)=\la x(u,t),y(u,t),0 \ra$}
	\label{fig:53}
	\begin{minipage}{0.23\linewidth}
		\centering
		\includegraphics[width=\linewidth]{Fig/53.pdf}
	\end{minipage}
	
	\vspace{0.3em}
	\parbox{0.95\linewidth}{\centering\fontsize{8pt}{10pt}\selectfont\textbf{Fuente:} elaboración propia.}
\end{minipage}
\vskip0.2cm


Lo anterior permite definir la función
{\small
\[
{\bf r}_2(t,u)= \sqrt[4]{(1-u)^3} \left( 4 \cos(t)-\cos(4t),4\sin(t)-\sin(4t),0\right)+ \sqrt[4]{u^3} \left(1,0,3\right).
\]
}

Las figuras de las superficies generadas por las funciones 
${\bf r}_1(t,u)$ y ${\bf r}_2(t,u)$ se muestran con diferentes colores en la siguiente figura. 

\begin{figure}[H]
	\begin{center}
		\captionof{figure}{\small Superficies determinadas por ${\bf r}_1(u,t)$ y ${\bf r}_2(u,t)$.}
		\includegraphics[height=4.5cm]{Img/78c.pdf}
		\includegraphics[height=4.5cm]{Img/78a.pdf}
		\includegraphics[height=4.5cm]{Img/78b.pdf}
	\end{center}
	\parbox{0.95\linewidth}{\centering\fontsize{8pt}{10pt}\selectfont{\bf Fuente:} elaboración propia.}
\end{figure}

Las siguientes instrucciones permiten obtener estos gráficos.

\begin{tcolorbox}[breakable,enhanced,title=\bf Instrucción \verb"Mathematica"]
\begin{Verbatim}[formatcom=\letraverbatim]
p={4*Cos[t]-Cos[4*t],4*Sin[t]-Sin[4*t],0};
r=(1-u)^(3/4)*p+u^(2/4)*{1,1,3};
ParametricPlot3D[{r,p*u},{u,0,1},{t,0,2Pi},
PlotStyle->{Orange,Cyan}]
\end{Verbatim}
\end{tcolorbox}

La función generada es ${\bf r}(u,t)= \langle(2\sqrt{u}+1-u^2)\cos t, (2\sqrt{u}+1-u^2) \sin t,\sqrt{u}\rangle,$
$ t\in [0,2\pi],\, u\in[0,1].$ 

\vskip0.2cm
\vspace{0pt}
\textcolor{black}{\scriptsize\ding{110}}\quad Las circunferencias parametrizadas por
$\ve{\alpha}(t)= \langle \cos(t),\sin(t),0 \rangle $ y $\ve{\beta}(t)= \langle 2\cos(t),2\sin(t),1 \rangle $,
se unen utilizando las funciones $g(u)= \sqrt{u}$ y $f(u)=1-u^2$; mediante la combinación
${\bf r}(u,t)= f(u)\ve{\alpha}(t) +g(u)\ve{\beta}(t)= \langle (2-2u^2+\sqrt{u})\cos t,(2-2u^2+\sqrt{u})\sin t,1- u^2\rangle$
con $t\in [0,2\pi],\,u\in[0,1].$


\vskip0.2cm
\noindent
\begin{minipage}{\textwidth}
	\centering
	\captionof{figure}{Circunferencias conectadas}
	\label{Circunferencias_conectadas}
	\begin{minipage}{0.32\linewidth}
		\centering
		\includegraphics[width=\linewidth]{Fig/54.pdf}
	\end{minipage}
	
	\vspace{0.3em}
	\parbox{0.95\linewidth}{\centering\fontsize{8pt}{10pt}\selectfont\textbf{Fuente:} elaboración propia.}
\end{minipage}
\vskip0.2cm


\textcolor{black}{\scriptsize\ding{110}}\quad La curva parametrizada por las ecuaciones $x(t)= 16\sin ^{3}(t),$ 
$y(t)=13\cos (t) -5\cos \left( 2t\right) -2\cos \left( 3t\right) -\cos \left(4t\right),$
con $t\in \lbrack 0,2\pi ]$, gira alrededor del eje vertical y genera una superficie.
Al hacer coincidir el eje vertical con el eje $z,$ entonces la función que describe la superficie en mención, para $t\in \lbrack 0,2\pi ],$ $u\in \lbrack 0,1],$ es
\[
\resizebox{\textwidth}{!}{$
	\mathbf{r}(t,u) = \langle 16\sin^3(t)\cos u,\ 16\sin^3(t)\sin u,\ 13\cos(t) - 5\cos(2t) - 2\cos(3t) - \cos(4t) \rangle
	$}
\]

\vspace{-4mm}
\begin{figure}[H]
	\centering
	% 1. Reducimos el espacio que el caption deja debajo de sí mismo
	\captionsetup{skip=2pt} 
	\caption{Superficies parametrizadas por $a(t, u)$ y $b(t, u)$}
	\label{fig:superficies}
	
	% 2. Eliminamos los vspace negativos entre imágenes si están en la misma línea
	% y usamos un vspace positivo pequeño para separarlas del caption
	\vspace{2pt} 
	\includegraphics[width=0.21\linewidth]{Img/49a.pdf} \hspace{1em}
	\includegraphics[width=0.21\linewidth]{Img/49b.pdf} \hspace{1em}
	\includegraphics[width=0.21\linewidth]{Fig/55.pdf}
	
	% 3. Controlamos el espacio de la fuente de manera independiente
	\vspace{0.2cm} 
	\parbox{0.95\linewidth}{\centering\fontsize{8pt}{10pt}\selectfont\textbf{Fuente:} elaboración propia.}
\end{figure}


Esta figura se genera ejecutando las siguientes instrucciones:

\begin{tcolorbox}[breakable,enhanced,title=\bf Instrucción \verb"Mathematica"] 
\begin{Verbatim}[formatcom=\letraverbatim]
ParametricPlot[{16*Sin[t]^3,13*Cos[t]-5*Cos[2*t]
-2*Cos[3*t]-Cos[4*t]},{t,0,2Pi}];
ParametricPlot3D[{16*Sin[t]^3*Sin[u],(16*Sin[t]^3)
*Cos[u],13*Cos[t]-5*Cos[2*t]-2*Cos[3*t]-Cos[4*t]},
{t,0,2Pi},{u,0,2*Pi},PlotPoints->80]
\end{Verbatim}
\end{tcolorbox}

\vspace{-0.2em}
Presentamos enseguida otros ejemplos de superficies parametrizadas bajo ciertas condiciones.
\subsection*{Superficie acotada por dos paraboloides} \index{Paraboloide} 
Las ecuaciones $8y=9-5x^{2}-6z^{2}$ y $16y=4z^{2}-x^{2},$
representan gráficamente un paraboloide y un paraboloide hiperbólico, al 
eliminar $y$ de estas expresiones se obtiene la ecuación $9x^{2}+16z^{2}=18,$ que corresponde a una elipse en el plano $xz.$ La región interior a esta curva se parametriza como 
$x=\sqrt{2}u\cos (t) ,$ $z=\frac{3\sqrt{2}}{4}u\sin (t) ,$ para $t\in \lbrack 0,2\pi ],$ $u\in \lbrack 0,1].$ Al reemplazar estas expresiones en las ecuaciones del paraboloide y del paraboloide hiperbólico, se obtiene:
\[
8y=9-5x^{2}-6z^{2}= 9-\frac{13}{4}u^{2}\cos ^{2}(t) -\frac{27}{4}u^{2},
\]
entonces
\[
y= \frac{9}{8}-\frac{13}{32}u^{2}\cos ^{2}(t) -\frac{27}{32}u^{2}.
\]
Por otra parte,
\[
16y=4z^{2}-x^{2}= \frac{9}{2}u^{2}-\frac{13}{2}u^{2}\cos^{2}(t) ,
\]
luego
\[
y= \frac{9}{32}u^{2}-\frac{13}{32}u^{2}\cos ^{2}(t) .
\]
Las superficies señaladas y el sólido  en mención se muestran enseguida.

\begin{figure}[H]
	\centering
	% 1. Centramos la leyenda respecto a todo el ancho de la página
	\caption{Superficies parametrizadas por $a(t,u)$ y $b(t,u)$}
	\label{fig:superficies_parametrizadas}
	\vspace{2pt} % Espacio pequeño para acercar la leyenda a las figuras
	
	% 2. Reducimos el tamaño de las minipages (de 0.32 a 0.28) para disminuir las figuras
	\begin{minipage}[b]{0.28\linewidth}
		\centering
		% Usamos scale o una altura menor (ej. 3.5cm) para homogeneizar
		\includegraphics[width=\linewidth, height=3.0cm, keepaspectratio]{Img/71b.pdf}
	\end{minipage}
	\hfill
	\begin{minipage}[b]{0.28\linewidth}
		\centering
		\includegraphics[width=\linewidth, height=3.0cm, keepaspectratio]{Img/71d.pdf}
	\end{minipage}
	\hfill
	\begin{minipage}[b]{0.28\linewidth}
		\centering
		\includegraphics[width=\linewidth, height=3.5cm, keepaspectratio]{Fig/56.pdf}
	\end{minipage}
	
	% 3. Espacio para la fuente
	\vspace{0.4cm}
	\parbox{0.95\linewidth}{\centering\fontsize{8pt}{10pt}\selectfont\textbf{Fuente:} elaboración propia.}
\end{figure}

\vskip0.3cm
Con lo señalado antes, la superficie sobre el paraboloide acotada por el
paraboloide hiperbólico se parametriza como
{\small
\[
\mathbf{a}(t,u) = \left\langle \sqrt{2}u\cos(t),\,
\frac{9}{8} - \frac{13}{32}u^{2}\cos^{2}(t) - \frac{27}{32}u^{2},\,
\frac{3\sqrt{2}}{4}u\sin(t) \right\rangle,
\]
}
mientras que la superficie sobre el paraboloide hiperbólico acotada por el paraboloide está dada por
{\small
\[
\mathbf{b}(t,u) = \left\langle \sqrt{2}u\cos(t),\,
\frac{9}{32}u^{2} - \frac{13}{32}u^{2}\cos^{2}(t),\,
\frac{3\sqrt{2}}{4}u\sin(t) \right\rangle.
\]
}
En ambos casos, $t \in [0, 2\pi]$ y $u \in [0, 1]$. 
Los gráficos se obtienen al ejecutar las siguientes instrucciones. Al ingresar el código se usó la identidad $2\cos^2(t) = 1 + \cos(2t)$.

\begin{tcolorbox}[breakable,enhanced,title=\bf Instrucción \verb"Mathematica"] 
\begin{Verbatim}[formatcom=\letraverbatim]
ParametricPlot3D[{{Sqrt[2]*u*Cos[t],(72-67*u^2
-13*u^2*Cos[2*t])/64,3*Sqrt[2]/4*u*Sin[t]},
{Sqrt[2]*u*Cos[t],u^2*(5-13*Cos[2*t])/64,
3*Sqrt[2]/4*u*Sin[t]}},{t,0,2*Pi},{u,0,1},
PlotStyle->{Orange,Cyan},PlotPoints->80]
\end{Verbatim}
\end{tcolorbox}

\subsection*{Superficies creadas por una cardioide y una elipse}
La cardioide con ecuación $r=1-\cos(t)$, se puede parametrizar,  para $t\in [0,2\pi]$ como 
$\langle(1-\cos(t))\cos(t),(1-\cos(t))\sin(t)\rangle$ y la elipse parametri- \\ zada en la forma $\langle 3\cos(t),2\sin(t)\rangle,$ $t\in [0,2\pi]$, determinan una región en el plano acotada por sus gráficos.

\vskip0.2cm
\noindent
\begin{minipage}{\textwidth}
	\centering
	\captionof{figure}{Superficie sobre el cilindro parabólico $z=1+(1-y)^2/4$}
	\label{Superficie_cilindro_parabólico}
	\begin{minipage}{0.37\linewidth}
		\centering
		\vspace{-3mm}
		\includegraphics[width=\linewidth]{Fig/57.pdf}
	\end{minipage}
	\hspace{1em}
	\begin{minipage}{0.37\linewidth}
		\centering
		\vspace{-3mm}
		\includegraphics[width=\linewidth]{Fig/58.pdf}
	\end{minipage}
	
	\vspace{0.3em}
	\parbox{0.95\linewidth}{\centering\fontsize{8pt}{10pt}\selectfont\textbf{Fuente:} elaboración propia.}
\end{minipage}
\vskip0.2cm

Una parametrización de la región del plano entre la elipse y el cardioide, se determina
formando una combinación convexa entre dos puntos, uno en cada curva, es decir
$(1-u)\langle(1-\cos(t))\cos(t),(1-\cos(t))\sin(t)\rangle
+u\langle 3\cos(t),2\sin(t)\rangle,$
esta expresión se simplifica en la forma 
{\small $${\bf r}(u,t)=\big \langle((u-1)\cos(t)+1+2u)\cos(t), (1+u+(u-1)\cos(t))\sin(t)\big \rangle,$$}
con $t\in [0,2\pi],\, u\in [0,1]$.
Esta región la proyectaremos sobre diferentes superficies para obtener su respectiva
parametrización. 
\vskip0.3cm
\subsubsection*{Superficie sobre el cilindro con ecuación \texorpdfstring{$z=1+(1-y)^2/4$}{z=1+(1-y)²/4}}

\vskip0.3cm
La ecuación $z=1+(1-y)^2/4$ nos permite llevar esta región del plano al espacio. 
Con lo mostrado anteriormente, solo nos falta determinar la coordenada $z$, esto es
$z=1+\frac{1}{4}(1-y)^2=1+\frac{1}{4}(1-((1+u+(u-1)\cos(t))\sin(t)))^2,$
con lo cual, la función que representa la parte del cilindro sobre la región que se considera es:
{\small
\begin{multline*}
{\bf r}(u,t)=\big \langle((u-1)\cos(t)+1+2u)\cos(t), (1+u+(u-1)\cos(t))\sin(t),\\
1+\frac{1}{4}(1-((1+u+(u-1)\cos(t))\sin(t)))^2 \big \rangle, 
\end{multline*}
}
con $t\in [0,2\pi],\, u\in [0,1].$ Una manera distinta de ver la misma superficie es considerarla como la parte del 
cilindro con ecuación $z=1+(1-y)^2/4$, que se encuentra entre el cilindro cardioidico
\index{Cilindro cardiodico}
con ecuación polar $r=1-\cos(t)$ y el cilindro elíptico con ecuación 
$\frac{x^2}{9}+\frac{y^2}{4}=1.$ La figura de esta situación se muestra a continuación.
\vspace{-1mm}
\vskip0.2cm
\noindent
\begin{minipage}{\textwidth}
	\centering
	\captionof{figure}{\small Superficie sobre el cilindro parabólico $z=1+(1-y)^2/4$}
	\begin{minipage}{0.24\linewidth}
		\centering
		\includegraphics[width=\linewidth]{Img/43b.pdf}
	\end{minipage}
	\hspace{1em}
	\begin{minipage}{0.24\linewidth}
		\centering
		\includegraphics[width=\linewidth]{Img/43a.pdf}
	\end{minipage}
	
	\vspace{0.3em}
	\parbox{0.95\linewidth}{\centering\fontsize{8pt}{10pt}\selectfont\textbf{Fuente:} elaboración propia.}
\end{minipage}

\vskip0.2cm
Estos gráficos se elaboraron en \verb"Mathematica".
En el plano $xy$ se forma una combinación convexa entre la función que determina la elipse y la cardioide.
Esta región se proyecta sobre el paraboloide $z=1+(y-1)^2/4$, la cual se grafica en el espacio.
Los cilindros se grafican por separado, definiendolos como \verb"a" y \verb"b",
que se muestran en un solo gráfico.

\begin{tcolorbox}[breakable,enhanced,title=\bf Instrucción \verb"Mathematica"] 
\begin{Verbatim}[formatcom=\letraverbatim]
{x1,y1}={(1-Cos{t})},*Cos{t},(1-Cos{t})*Sin{t});
{x2,y2}={3*Cos[t],2*Sin[t]};
{x,y}={1-u}*{x1,y1}+u*{x2,y2};
ParametricPlot[{x,y},{t,0,2*Pi},{u,0,1}]
ParametricPlot3D[{x,y,1+(1-y)^2/4},{u,0,1}{t,0,2*Pi},
PlotRange->{{-3, 3}, {-2, 2}},
MeshStyle->Gray]
a=ContourPlot3D[{z==1+(1-y)^2/4},
{x,-4,4},{y,-4,4},{z,-1,5},
ContourStyle->RGBColor[0.1,0.4,0.3], Mesh->False];
b=ParametricPlot3D[{{x1,y1,u},{x2,y2, u}},{u,0,5},
{t,0,2*Pi}, 
Mesh->False];
Show[a, b];
	\end{Verbatim}
\end{tcolorbox}

\subsubsection*{Superficie sobre el cilindro \texorpdfstring{$z=\cos (2x)$}{z=cos(2x)}}
En este caso, la parametrización de la superficie es
{\small
\begin{multline*}
{\bf r}\left( u,t\right) = \big\langle \cos (t) \left( \cos (t)
\left( u-1\right) +1+2u\right) ,\sin (t) ( 1+u+\\ 
(u-1) \cos (t)) , \cos \big( 2\cos (t)\left( \cos (t) \left( u-1\right) +1+2u\right) \big) \big \rangle,
\end{multline*}
}
con $t \in [0,2\pi]$, $u \in [0,1]$, dicha superficie se muestra a continuación.
\vskip0.2cm
\noindent
\begin{minipage}{\textwidth}
	\centering
	\captionof{figure}{\small Superficie sobre el cilindro $z=\cos (2x)$}
	\begin{minipage}{0.24\linewidth}
		\centering
		\includegraphics[width=\linewidth]{Img/46a.pdf}
	\end{minipage}
	\hspace{1em}
	\begin{minipage}{0.24\linewidth}
		\centering
		\includegraphics[width=\linewidth]{Img/46b.pdf}
	\end{minipage}
	
	\vspace{0.3em}
	\parbox{0.95\linewidth}{\centering\fontsize{8pt}{10pt}\selectfont\textbf{Nota:} elaboración propia.}
\end{minipage}

\vspace{0.3cm}
Las instrucciones para conseguirlo son las siguientes:
primero, se define la función ${\bf r}$, que describe la región entre la elipse y la cardioide;
luego, se define la función que parametriza la superficie y 
se muestra su gráfico. 

\begin{tcolorbox}[breakable,enhanced,title=\bf Instrucción \verb"Mathematica"] 
\begin{Verbatim}[formatcom=\letraverbatim]
r=(1-u)*{(1-Cos[t])*Cos[t],(1-Cos[t])*Sin[t]}+
u*{3*Cos[t],2*Sin[t]};
ParametricPlot3D[{r[[1]],r[[2]],Cos[2*r[[1]]]},
{u,0,1},{t,0,2*Pi},PlotRange->{{-3,3},{-3,3},{-1,1}},
PlotStyle->RGBColor[0.8,0.2,0.3],Mesh->False,
PlotPoints->90,Ticks->{{-2,0,2},{-2,0,2},{-1,0,1}}]
\end{Verbatim}
\end{tcolorbox}

\subsubsection*{Superficie sobre el cilindro \texorpdfstring{$z=\sin (2x)$}{z=sin(2x)}}
Sobre este cilindro la parametrización que se utiliza es:
{\small
\begin{multline*}
{\bf r}(u,t) = \big\langle \cos (t) \left( \cos (t)\left( u-1\right) +1+2u\right),
\sin(t)(1+u+\\ (u-1)\cos(t)), \sin \big( 2\cos (t)\left( \cos (t) \left( u-1\right) +1+2u\right)\big) \big\rangle,
\end{multline*}
}
en donde $t \in [0,2\pi]$, $u \in [0,1]$.
Se presenta ahora la figura y las instrucciones
que se requieren para obtenerlo:


\vskip0.2cm
\noindent
\begin{minipage}{\textwidth}
	\centering
	\captionof{figure}{\small Superficie sobre el cilindro $z=\sin(2x)$}
	\begin{minipage}{0.23\linewidth}
		\centering
		\includegraphics[width=\linewidth]{Img/47b.pdf}
	\end{minipage}
	\begin{minipage}{0.23\linewidth}
		\centering
		\includegraphics[width=\linewidth]{Img/47a.pdf}
	\end{minipage}
	
	\vspace{0.3em}
	\parbox{0.95\linewidth}{\centering\fontsize{8pt}{10pt}\selectfont\textbf{Fuente:} elaboración propia.}
\end{minipage}
\vskip0.2cm



\begin{tcolorbox}[breakable,enhanced,title=\bf Instrucción \verb"Mathematica"]  
\begin{Verbatim}[formatcom=\letraverbatim]
r=(1-u)*{(1-Cos[t])*Cos[t],(1-Cos[t])*Sin[t]}+
u*{3*Cos[t],2*Sin[t]};
p={r[[1]],r[[2]],Sin[2*r[[1]]]};
ParametricPlot3D[p,{u,0,1},{t,0,2*Pi},
PlotRange->{{-3,3},{-3,3},{-1,1}},Mesh->False,
PlotStyle->RGBColor[0.8,0.2,0.3],PlotPoints->80]
\end{Verbatim}
\end{tcolorbox}

\subsubsection*{Superficie sobre el cilindro \texorpdfstring{$z=\frac{1}{5}y^{2}-\frac{1}{4}(x-1) ^{2}$}{z=(1/5)y²-(1/4)(x-1)²}}

La superficie mencionada se presenta a continuación.

\vskip0.2cm
\noindent
\begin{minipage}{\textwidth}
	\centering
	\captionof{figure}{\small Superficie sobre el cilindro $z=\frac{1}{5}y^{2}-\frac{1}{4}(x-1) ^{2}$}
	\begin{minipage}{0.24\linewidth}
		\centering
		\includegraphics[width=\linewidth]{Img/45a.pdf}
	\end{minipage}
	\begin{minipage}{0.24\linewidth}
		\centering
		\includegraphics[width=\linewidth]{Img/45b.pdf}
	\end{minipage}
	
	\vspace{0.3em}
	\parbox{0.95\linewidth}{\centering\fontsize{8pt}{10pt}\selectfont\textbf{Fuente:} elaboración propia.}
\end{minipage}
\vskip0.2cm


Como se señaló antes, nos falta determinar la coordenada $z$, esto es 
{\footnotesize
	\begin{align*}
z&=\frac{1}{5}y^{2}-\frac{1}{4}\left( x-1\right) ^{2} \\	&=\frac{1}{5}\big( \sin(t) \left( 1+u+(u-1) \cos(t) \right) \big)^{2}-\frac{1}{4}\big(\cos(t)\left(\cos(t) \left( u-1\right) +1+2u\right)-1\big)^{2},
	\end{align*}
}
entonces la parametrización de la superficie es
{\small
\begin{multline*}
{\bf r}(u,t) = \bigg\langle \cos (t) \left( \cos (t)
\left( u-1\right) +1+2u\right) ,\sin (t)(1+u
+(u-1) \cos (t)) , \\ \frac{1}{5}\big( \sin(t) \left( 1+u+(u-1) \cos(t) \right)\big) ^{2} \\
-\frac{1}{4}\big( \cos (t) \left( \cos (t) \left( u-1\right) +1+2u\right) -1\big) ^{2} \bigg \rangle. 
\end{multline*}
}
Para $t \in [0,2\pi]$, $u \in [0,1]$, las instrucciones 
para conseguir el gráfico se muestran enseguida.

%\vspace{-4mm}
\begin{tcolorbox}[breakable,enhanced,title=\bf Instrucción \verb"Mathematica"] 
\begin{Verbatim}[formatcom=\letraverbatim]
r=(1-u)*{(1-Cos[t])*Cos[t],(1-Cos[t])*Sin[t]}+
u*{3*Cos[t],2*Sin[t]};
p={r[[1]],r[[2]],-1/4*(r[[1]]-1)^2+1/5*(r[[2]])^2};
ParametricPlot3D[p,{u,0,1},{t,0,2*Pi},
PlotRange->{{-3,3},{-2,2},{-4.2,1}},Mesh->False,
PlotStyle->RGBColor[0.2,0.5,0.5],PlotPoints->80]
\end{Verbatim}
\end{tcolorbox}

\subsection*{Superficies acotadas por cilindros elípticos}

Las expresiones $4x^2 + y^2 = 36$ y $x^2 + 4y^2 = 4$ representan dos cilindros elípticos con eje $z$ como eje de simetría. Sin embargo, se puede dar una mejor idea de sus ecuaciones representando dos elipses que pueden parametrizarse respectivamente como $\vec{a}(t) = \langle 3\cos(t),\, 6\sin(t) \rangle$ y $\vec{b}(t) = \langle 2\cos(t),\, \sin(t) \rangle$.




\vskip0.2cm
\noindent
\begin{minipage}{\textwidth}
	\centering
	\captionof{figure}{Región parametrizada por ${\bf r}(u,t)=\left\langle \cos(t)(3-u),\, \sin(t)(6-5u) \right\rangle$}
	\label{fig:59}
	\begin{minipage}{0.23\linewidth}
		\centering
		\includegraphics[width=\linewidth]{Img/1a.pdf}
	\end{minipage}
	\begin{minipage}{0.28\linewidth}
		\centering
		\includegraphics[width=\linewidth]{Fig/59.pdf}
	\end{minipage}
	\begin{minipage}{0.28\linewidth}
		\centering
		\includegraphics[width=\linewidth]{Fig/60.pdf}
	\end{minipage}
	
	\vspace{0.3em}
	\parbox{0.95\linewidth}{\centering\fontsize{8pt}{10pt}\selectfont\textbf{Fuente:} elaboración propia.}
\end{minipage}
\vskip0.2cm


Con dos puntos $P$ y $Q$, simplificando la expresión: 
\begin{align}
	(1 - u)P + uQ &= (1 - u)(3\cos(t),\, 6\sin(t)) + u(2\cos(t),\, \sin(t)) \notag \\
	&= (\cos(t)(3 - u),\, \sin(t)(6 - 5u)). \tag{3.7} \label{1a}
\end{align}
Para $u \in [0,1]$, esta función parametriza la región que se halla entre estas dos curvas. La región en mención, que denominaremos $R$, se proyectará sobre un plano y posteriormente sobre un paraboloide; se presentará su construcción con sus respectivas parametrizaciones.

\subsubsection*{Superficie sobre un plano}
Se puede representar la porción del plano con ecuación $6z=2x+y+12$ sobre la región acotada  por los dos cilindros. 

\begin{figure}[H]
	\centering
	\captionof{figure}{\small Cilindros elípticos intersecados por el plano $6z=2x+y+12$}
	\label{fig:cilindros_intersecados}
	\vspace{2pt} % Acerca la leyenda a las imágenes
	
	% Reducción uniforme: width de 4cm a 3.5cm o 0.28\linewidth
	\includegraphics[width=0.28\linewidth, height=3cm, keepaspectratio]{Img/1c.pdf}
	\hspace{1cm} % Espacio horizontal entre ambas
	\includegraphics[width=0.28\linewidth, height=3cm, keepaspectratio]{Img/1d.pdf}
	
	\vspace{0.4cm} % Separación con la fuente
	\parbox{0.95\linewidth}{\centering\fontsize{8pt}{10pt}\selectfont\textbf{Fuente:} elaboración propia.}
\end{figure}

Para determinar $z$ se puede proceder como sigue
\vspace*{-4mm} 
\[
\begin{aligned}
	z &= \frac{1}{6}\left( 2\cos (t) \left( 3-u\right) +\sin (t) \left( 6-5u\right) +12\right) \\ 
	&= \left( 1-\frac{1}{3}u\right)\cos (t) +\left( 1-\frac{5}{6}u\right) \sin (t) +2
\end{aligned}
\]
Por último, la región $R$ sobre el plano se parametriza mediante la función:
\vspace*{-2mm} 
{\footnotesize
	\[
	\overrightarrow{r}(u,t) =\big\langle \cos (t) \left( 3-u\right) ,\sin (t) \left( 6-5u\right) ,
	\left( 1 - \! \frac{1}{3}u\right) \cos (t) + \left( 1 - \! \frac{5}{6}u\right) \sin (t) +2\big\rangle ,
	\]
}
\vspace*{-1mm} 
con $t\in \lbrack 0,2\pi ],\,\, u\in \lbrack 0,1].$


\vskip0.2cm
\noindent
\begin{minipage}{\textwidth}
	\centering
	\captionof{figure}{Región sobre el plano $6z=2x+y+12$ acotada por los cilindros elípticos}
	\label{Región_cilindros_elípticos}
	\begin{minipage}{0.50\linewidth}
		\centering
		\includegraphics[width=\linewidth]{Fig/61.pdf}
	\end{minipage}
	
	\vspace{0.3em}
	\parbox{0.95\linewidth}{\centering\fontsize{8pt}{10pt}\selectfont\textbf{Fuente:} elaboración propia.}
\end{minipage}
\vskip0.2cm



\subsubsection*{Superficie sobre un paraboloide} \index{Paraboloide} 
Con ayuda de la expresión (\ref{1a}), se puede parametrizar la parte del
parabo- \\ loide con ecuación $z=6-\left( x-1\right) ^2-y^2,$ cuya proyección 
en el plano $xy$ es la región $R,$ para este fin determinaremos $z$, esto es
\begin{eqnarray*} z&=&  6-\left( \left( 3-u\right)\cos(t)-1\right)^2-\left(
\left( 6-5u\right) \sin (t) \right) ^2 \\
&=& -31+\left( 27-54u+24u^{2}\right) \cos ^{2}(t) +\left( 6-2u\right) \cos (t) +60u-25u^{2}
\end{eqnarray*}


\vskip0.2cm
\noindent
\begin{minipage}{\textwidth}
	\centering
	\captionof{figure}{Superficie sobre el paraboloide acotada por los cilindros elípticos}
	\label{Superficie_paraboloide_acotada_cilindros_elípticos}
	\begin{minipage}{0.23\linewidth}
		\centering
		\includegraphics[width=\linewidth]{Img/1e.pdf}
	\end{minipage}
	\begin{minipage}{0.23\linewidth}
		\centering
		\includegraphics[width=\linewidth]{Img/1f.pdf}
	\end{minipage}
	\begin{minipage}{0.28\linewidth}
		\centering
		\includegraphics[width=\linewidth]{Fig/62.pdf}
	\end{minipage}
	
	\vspace{0.3em}
	\parbox{0.95\linewidth}{\centering\fontsize{8pt}{10pt}\selectfont\textbf{Fuente:} elaboración propia.}
\end{minipage}
\vskip0.2cm


Por lo tanto, la función 
\vspace*{-4mm} % Ajuste vertical después del párrafo
\begin{multline*}
	\overrightarrow{r}(u,t) = \big\langle \cos(t)(3-u) ,\sin(t)( 6-5u) , \\
	-31+\left( 27-54u+24u^{2}\right) \cos ^{2}(t) +\left( 6-2u\right) \cos (t) +60u-25u^{2} \big\rangle
\end{multline*}

\noindent con $t\in \lbrack 0,2\pi ],$ $u\in \lbrack 0,1],$
con $ t\in \lbrack 0,2\pi ],$ $u\in \lbrack 0,1],$  permite representar gráficamente la situación mencionada.
% Las curvas de intersección de los cilindros elípticos y el paraboloide se determinan reemplazando en el valor 
% de $z,$ que para los dos casos es:
% $$z=6-\left( 3\cos (t) -1\right) ^{2}-\left( 6\sin (t) \right) ^{2}= -31+27\cos ^{2}(t) +6\cos (t) $$
% $$z=6-\left( 2\cos (t) -1\right) ^{2}-\left( \sin (t) \right) ^{2}= 4-3\cos ^{2}(t) +4\cos (t) $$
Las figuras de esta sección pueden obtenerse ejecutando las siguientes instrucciones:

\begin{tcolorbox}[breakable,enhanced,title=\bf Instrucción \verb"Mathematica"] 
\begin{Verbatim}[formatcom=\letraverbatim]
ContourPlot3D[{4*x^2+y^2==36, x^2+4*y^2==4,6*z==-2*x+y+12},
{x,-6,6},{y,-6,6},{z,0,4},PlotPoints->50]
ContourPlot3D[{6*z==-2*x+y+12, 4*x^2+y^2==36,x^2+4*y^2==4},
{x,-4,4},{y,-7,7},{z,-0.5,4.5},PlotPoints->80,
ContourStyle->{Gray,Yellow,Red}]
ContourPlot3D[{z==6-(x-1)^2-y^2,4*x^2+y^2==36,
x^2+4*y^2==4},
{x,-6,10},{y,-7,7},{z,-35,6},
ContourStyle->{Red,Yellow,Green}]
\end{Verbatim}
\end{tcolorbox}


\begin{tcolorbox}[breakable,enhanced,title=\bf Instrucción \verb"Mathematica"] 
\begin{Verbatim}[formatcom=\letraverbatim]
ParametricPlot3D[{{1 + 3*Cos[t],2*Sin[t],Sqrt[5*(Sin[t])^2 
-6*Cos[t]+6]},{1+3*Cos[t],2*Sin[t],-Sqrt[5*(Sin[t])^2
-6*Cos[t]+6]}},{t,0,2*Pi}, PlotStyle->{Directive[Black,
Thickness[0.015]],Directive[Black,Thickness[0.015]]}]
\end{Verbatim}
\end{tcolorbox}

\subsection*{Dos paraboloides} \index{Paraboloide!hiperbólico}
Las ecuaciones $4y=z^{2}-x^{2}\,$ y $\, y=4-x^{2}-z^{2}$ 
representan un paraboloide y un paraboloide hiperbólico. Estas figuras se muestran enseguida.

\vskip0.2cm
\noindent
\begin{minipage}{\textwidth}
	\centering
	\captionof{figure}{\small $4y=z^{2}-x^{2}$ y $y=4-x^{2}-z^{2}$}
	\begin{minipage}{0.24\linewidth}
		\centering
		\includegraphics[width=\linewidth]{Img/4a.pdf}
	\end{minipage}
	\hspace{1em} 
	\begin{minipage}{0.24\linewidth}
		\centering
		\includegraphics[width=\linewidth]{Img/4b.pdf}
	\end{minipage}
	
	\vspace{0.3em}
	\parbox{0.95\linewidth}{\centering\fontsize{8pt}{10pt}\selectfont\textbf{Fuente:} elaboración propia.}
\end{minipage}
\vskip0.2cm

Estas superficies se intersecan 
en una curva con ecuación cartesiana 
$3x^{2}+5z^{2}=16,$ que equivale a 
$\frac{x^{2}}{16/3}+\frac{z^{2}}{16/5}=1,$ la cual representa una elipse en
el plano $xz.$ 


\vskip0.2cm
\noindent
\begin{minipage}{\textwidth}
	\centering
	\captionof{figure}{$3x^{2}+5z^{2}\leq 16$}
	\label{fig:63}
	\begin{minipage}{0.34\linewidth}
		\centering
		\vspace{-4mm}
		\includegraphics[width=\linewidth]{Fig/63.pdf}
	\end{minipage}
	
	\vspace{4mm}
	\parbox{0.95\linewidth}{\centering\fontsize{8pt}{10pt}\selectfont\textbf{Fuente:} elaboración propia.}
\end{minipage}
\vskip0.2cm




La frontera puede parametrizarse en la forma
$x=\frac{4}{\sqrt{3}}\cos (t) ,\,\, z=\frac{4}{\sqrt{5}}\sin(t),$  $t\in \lbrack 0,2\pi ].$
Mientras que el interior se parametriza en la forma  
$x=\frac{4u}{\sqrt{3}}\cos (t) ,$ $z=\frac{4u}{\sqrt{5}}\sin(t),$ 
$t\in \lbrack 0,2\pi ]$ y $u\in \lbrack 0,1].$ Reemplazando $x$ y $z$ en la ecuación del paraboloide se sigue que
{\footnotesize
	\[
	y = 4-x^{2}-z^{2}= 4-\left( \frac{4}{\sqrt{3}}\cos \left( t\right) \right) ^{2}-\left( \frac{4}{\sqrt{5}}\sin \left( t\right) \right) ^{2} =  \frac{4}{5}-\frac{32}{15}\cos ^{2}\left( t\right)
	\]
}
Con lo anterior, se establece que la función vectorial:
\[
\overrightarrow{r}(t)=\left\langle \frac{4}{\sqrt{3}}\cos (t) ,\frac{4}{5}-\frac{32}{15}\cos ^{2}(t) ,\frac{4}{\sqrt{5}}\sin (t) \right\rangle,
\]
\noindent con $t\in \lbrack 0,2\pi ],$ representa la curva $C$, que es la intersección.



\noindent \textcolor{black}{\scriptsize\ding{110}}\quad La superficie sobre el paraboloide hiperbólico acotado por la curva $C$, se puede 
parametrizar reemplazando $x$ y $z$, para obtener el valor de $y$, 
\[
y=\frac{1}{4}\left(\left(\frac{4u}{\sqrt{5}}\sin (t) \right) ^{2}-\left( \frac{4u}{\sqrt{3}}\cos (t) \right) ^{2}\right) = \frac{4}{5}u^{2}-\frac{32}{15}u^{2}\cos ^{2}(t),
\]
	
por lo tanto, 
\[
\overrightarrow{r}(t,u)=\left\langle \frac{4u}{\sqrt{3}}\cos (t) , \frac{4}{5}u^{2}-\frac{32}{15}u^{2}\cos ^{2}(t) , \frac{4u}{\sqrt{5}}\sin (t)  \right\rangle,
\]
con $t\in [0,2\pi ],$ $u\in [ 0,1 ], $ es una manera de parametrizar la mencionada superficie.
	

%\vskip0.2cm
\noindent
\begin{minipage}{\textwidth}
	\centering
	\captionof{figure}{Superficie determinada por $\mathbf{r}(t,u)$}
	\label{fig:64}
	\begin{minipage}{0.32\linewidth}
		\centering
		\vspace{-3mm}
		\includegraphics[width=\linewidth]{Fig/64.pdf}
	\end{minipage}
	
	\vspace{0.3em}
	\parbox{0.95\linewidth}{\centering\fontsize{8pt}{10pt}\selectfont\textbf{Fuente:} elaboración propia.}
\end{minipage}
\vskip0.2cm

\noindent
\textcolor{black}{\scriptsize\ding{110}}\quad Ahora se desea parametrizar la parte del paraboloide acotado por la curva $C$.
Para este propósito, el valor de $y$ se consigue sustituyendo en la expresión adecuada:

%\vspace*{2mm} 
{\footnotesize
	\[
	y = 4-\left( \frac{4u}{\sqrt{3}}\cos(t)\right) ^{2}-\left(\frac{4u}{\sqrt{5}}\sin (t)\right)^{2} = 4-\frac{32}{15}u^{2}\cos ^{2}(t) -\frac{16}{5}u^{2},
	\]
}
\vspace{2mm} 
por lo tanto,
{\footnotesize
	\[
	\overrightarrow{r}(t,u)=\left\langle \frac{4u}{\sqrt{3}}\cos (t) , 4-\frac{32}{15}u^{2}\cos ^{2}(t) -\frac{16}{5}u^{2} , \frac{4u}{\sqrt{5}}\sin (t) \right\rangle,
	\]
}
con $t\in \lbrack 0,2\pi ],\, u\in \lbrack 0,1 ], $ representa la superficie deseada.

\begin{figure}[H]
	\centering 
	\captionof{figure}{\small Superficie acotada por los paraboloides $4y=z^{2}-x^{2}$ y $y=4-x^{2}-z^{2}$}
	\includegraphics[height=4.5cm]{Img/4c.pdf} \hspace{0.5cm} 
	\includegraphics[height=4.5cm]{Img/4d.pdf} 
	\parbox{0.95\linewidth}{\centering\fontsize{8pt}{10pt}\selectfont\textbf{Fuente:} elaboración propia.}
\end{figure}

\noindent{Esta figura se obtiene ejecutando las siguientes instrucciones:}
\vspace{0.5em}

\begin{tcolorbox}[breakable,enhanced,title=\bf Instrucción \verb"Mathematica"]  
\begin{Verbatim}[formatcom=\letraverbatim]
a=Show[ContourPlot3D[{4*y==z^2-x^2,y==4-x^2-z^2},{x,-3,3},
{y,-3,4},{z,-3,3},Mesh->None,PlotPoints->90,
ContourStyle->{{Gray,Orange},Opacity[0.6]}],
ParametricPlot3D[{4/Sqrt[3]*Cos[t],4-(4/Sqrt[5]*Sin[t])^2
+(4/Sqrt[3]*Cos[t])^2,4/Sqrt[5]*Sin[t]},{t,0,2*Pi},
PlotPoints->80,PlotStyle->{Purple,Thickness->0.015}]];

b=ParametricPlot3D[{4*u/Sqrt[3]*Cos[t],
4-32/15*u^2*Cos[t]^2-
16*u^2/5,4*u/Sqrt[5]*Sin[t]},{t,0,2*Pi},{u,0,1},
PlotPoints->90,
Ticks->{{-2,0,2},{0,2,4},{-1,0,1}},PlotStyle->RGBColor[0.9,
0.1,0.1,1],MeshStyle->RGBColor[0.3,0.1,0.1,1]];

Show[a,b]
\end{Verbatim}
\end{tcolorbox}

\subsection*{Una esfera y un cilindro elíptico}

Consideraremos la esfera 
$x^2 + y^2 + z^2 = 16$ y el cilindro elíptico $\frac{(x - 1)^2}{9} +\frac{y^2}{4} = 1$. 
La región interior al cilindro se parametriza en la forma $x-1=3u\cos(t)$, $y=2u\sin(t)$, 
con $u \in[0,1]$ y $t \in [0,2\pi]$. Mientras que para la frontera usaremos $x=1+3\cos(t)$, $y=2\sin(t)$.
\vskip0.3cm

\begin{figure}[H]
	\begin{center}
		\centering
		\captionof{figure}{\small Esfera y un cilindro elíptico}
		\includegraphics[width=3.5cm,height=3.25cm]{Img/2a.pdf}\hspace{0.5cm}
		\includegraphics[width=3.5cm,height=3.25cm]{Img/2b.pdf}
		\parbox{0.95\linewidth}{\centering\fontsize{8pt}{10pt}\selectfont\textbf{Fuente:} elaboración propia.}
	\end{center}
\end{figure}

Con lo anterior, se obtienen las ecuaciones de la curva de intersección del cilindro y la esfera que se determinan reemplazando el valor de $z$ en la ecuación, es decir
$\left( 1+3\cos (t) \right) ^{2}+\left( 2\sin (t) \right) ^{2}+z^{2}=16$, de donde \linebreak
\[
z=\pm \sqrt{5\sin ^{2}(t) -6\cos (t) +6}.
\]
\vskip0.1cm 
Por lo tanto, una parametrización para las curvas en mención está dada por
{\small
\[
\vec{r}(t) = \left\langle 1 + 3\cos(t),\, 2\sin(t),\, \pm \sqrt{5\sin^2(t) - 6\cos(t) + 6} \right\rangle,
\quad t \in [0, 2\pi].
\]
}

Una parte de la curva determina la raíz positiva, mientras que la raíz negativa representa la otra parte de la curva. La figura de esta situación se muestra a continuación.


\vskip0.2cm
\noindent
\begin{minipage}{\textwidth}
	\centering
	\captionof{figure}{\small Esfera y un cilindro elíptico con su curva de intersección}
	\begin{minipage}{0.24\linewidth}
		\centering
		\includegraphics[width=\linewidth]{Img/2c.pdf}
	\end{minipage}
	\begin{minipage}{0.24\linewidth}
		\centering
		\includegraphics[width=\linewidth]{Img/2d.pdf}
	\end{minipage}
	
	\vspace{0.3em}
	\parbox{0.95\linewidth}{\centering\fontsize{8pt}{10pt}\selectfont\textbf{Fuente:} elaboración propia.}
\end{minipage}
\vskip0.2cm




Para la superficie de la esfera que se halla dentro del cilindro, el valor de $z$ está dado por:
\vspace*{-3mm} % Ajuste vertical para pegar la ecuación al párrafo
\[
\left( 1+3u\cos (t) \right) ^{2}+\left( 2u\sin (t) \right) ^{2}+z^{2}=16,
\]
\vspace{4mm} % Espacio vertical para separar la ecuación del texto siguiente
\noindent de donde, \[z= \pm \sqrt{15-6u\cos (t) -9u^{2}+5u^{2}\sin ^{2}(t)}\] con esto se establece que
$${\bf r}(u,t) =\left\langle 1+3u\cos(t),u\sin(t),\sqrt{15-6u\cos(t)-9u^{2}+5u^{2}\sin^{2}(t)}\right\rangle,$$

con $ t \in [0, 2\pi], \,\, u\in [0,1]  $, es una parametrización de la parte de la esfera dentro del cilindro.

\vskip0.2cm
\noindent
\begin{minipage}{\textwidth} % <--- MINIPAGE: Fija la figura al flujo de texto
	\centering % <--- \centering es más compacto que \begin{center}
		\captionof{figure}{\small Superficie sobre la esfera acotada por el cilindro elíptico}
		\includegraphics[width=2.8cm,height=2.90cm]{Img/2e.pdf}\hspace{1cm}
		\includegraphics[width=2.8cm,height=2.9cm]{Img/2f.pdf}
		\parbox{0.95\linewidth}{\centering\fontsize{8pt}{10pt}\selectfont\textbf{Fuente:} elaboración propia.}
	\end{minipage}

\vspace{0.3cm}
\noindent 
Para representar la parte del cilindro que está dentro de la esfera, podemos tomar dos puntos $P = \langle 1 + 3\cos(t),\, 2\sin(t),\, \sqrt{6 - 6\cos(t) + 5\sin^2(t)} \rangle$ y $Q = \langle 1 + 3\cos(t),\, 2\sin(t),\, -\sqrt{6 - 6\cos(t) + 5\sin^2(t)} \rangle$. Con ellos se parametriza el segmento rectilíneo que los une. Simplificando la expresión $(1 - u)P + uQ$, con $u \in [0, 1]$, se obtiene la función $\vec{r}(u, t)$ dada por:	
	\[
	\vec{r}(u, t) = \langle 1 + 3\cos(t),\, 2\sin(t),\, (1 - 2u)\sqrt{6 - 6\cos(t) + 5\sin^2(t)} \rangle
	\]	
Esta función, con $t \in [0, 2\pi]$ y $u \in [0, 1]$, representa una parametrización de la superficie sobre el cilindro que se encuentra dentro de la esfera.

\begin{figure}[H]
	\centering
	\caption{Superficie sobre el cilindro elíptico acotada por la esfera}
	\label{Superficie_cilindro_elíptico_acotada_esfera}
	\begin{minipage}[b]{0.24\linewidth}
		\centering
		\includegraphics[width=\linewidth]{Img/2g.pdf}
	\end{minipage}
	\hspace{1em}
	\begin{minipage}[b]{0.25\linewidth}
		\centering
		\includegraphics[width=\linewidth]{Img/2h.pdf}
	\end{minipage}
	\hspace{1em}
	\begin{minipage}[b]{0.24\linewidth}
		\centering
		\includegraphics[width=\linewidth]{Fig/65.pdf}
	\end{minipage}
	\parbox{0.95\linewidth}{\centering\fontsize{9pt}{11pt}\selectfont{\bf Fuente:} elaboración propia.}
\end{figure}


Los gráficos se obtienen con ayuda de las siguientes instrucciones:

\begin{tcolorbox}[breakable,enhanced,title=\bf Instrucción \verb"Mathematica"] 
\begin{Verbatim}[formatcom=\letraverbatim]
Show[{ContourPlot3D[{x^2+y^2+z^2==1,(x-1)^2/9+y^2/4==1},
{x,-4,4},{y,-4,4},{z,-4,4.0},ContourStyle->{Directive[Orange,
Opacity[0.15]],Directive[Gray,Opacity[0.15]]},Mesh->False],
ParametricPlot3D[{1+3*u*Cos[t],2*u*Sin[t],Sqrt[15-6*u*Cos[t]
-9*u^2+2+5*u^2*(Sin[t])^2]},{u,0,1}{t,0,2Pi},PlotStyle->Red]}]

Show[{ContourPlot3D[{x^2+y^2+z^2==1,(x-1)^2/9+y^2/4==1},
{x,-4,4},{y,-4,4},{z,-4,4.0},ContourStyle->{Directive[Orange,
Opacity[0.15]],Directive[Gray,Opacity[0.15]]},Mesh->False],
ParametricPlot3D[{1+3*Cos[t], 2*Sin[t],(1-2*u)Sqrt[6
-6*Cos[t]+5*(Sin[t])^2]},{u,0,1},{t,0,2Pi},PlotStyle->Red]}]
\end{Verbatim}
\end{tcolorbox}


\section{Sólidos} \label{p1} \index{Parametrización!de sólidos}
Con los elementos anteriores parametrizamos algunos sólidos.

\subsection*{Sólido acotado por dos paraboloides}

Los paraboloides elípticos con ecuaciones $z=3-2x^2-y^2$ y $z=x^2+2y^2,$ se intersectan en una región determinada por $x^2+y^2=1,$ cuyo interior se representa con $x=u \cos v,$ $y=u \sin v$ con $u\in [0,1]$ y $v\in [0,2\pi].$


%\vspace*{-0.5\baselineskip} 
La coordenada $z$ de dos puntos $P$ y $Q$ sobre los paraboloides, es
\begin{align*}
z_1 &= 3-2(u\cos v)^2-(u\sin v)^2 \\ &= \frac{1}{2}(6-3u^2-u^2 \cos(2v)) \\
z_2 &= (u\cos v)^2+2(u\sin v)^2 \\ &= \frac{u^2}{2}(3-\cos(2v))
\end{align*}


\vskip0.2cm
\noindent
\begin{minipage}{\textwidth}
	\centering
	\captionof{figure}{\small Sólido acotado por paraboloides}
	\label{fig:300}
	\begin{minipage}{0.35\linewidth}
		\centering
		\includegraphics[width=\linewidth]{Fig/300.pdf}
	\end{minipage}
	
	\vspace{0.3em}
	\parbox{0.95\linewidth}{\centering\fontsize{8pt}{10pt}\selectfont\textbf{Fuente:} elaboración propia.}
\end{minipage}
\vskip0.2cm


\noindent
La región se parametriza uniendo los puntos $P(u\cos v,u\sin v,\frac{u^2}{2}(3-\cos(2t)))$ y 
$Q(u\cos v,u\sin v,\frac{1}{2} (6-3u^2-u^2 \cos(2t))),$
mediante segmentos rectilíneos. Es decir, la expresión $(1-w)P+wQ$ se simplifica y permite definir la función
$${\bf r}(u,v,w)= \langle u\cos v,u\sin v, 3(1-w)+u^2(6w-3-\cos(2t))/2 \rangle ,$$
$u\in [0,1],$ $v\in [0,2\pi],$ $w\in [0,1].$  Una representación de estas superficies se obtiene con las siguientes instrucciones:

\begin{tcolorbox}[breakable,enhanced,title=\bf Instrucción \verb"Mathematica"] 
\begin{Verbatim}[formatcom=\letraverbatim]
A=ContourPlot3D[{z==3-2*x^2-y^2,z==x^2+2*y^2},
{x,-2,2},{y,-2,2},{z,0,3},Mesh->False,
ContourStyle->{Directive[Orange,Opacity[0.8]],
Directive[Cyan,Opacity[0.95]]},PlotPoints->80];
B=ParametricPlot3D[{Cos[t],Sin[t],Cos[t]^2+
2*Sin[t]^2},{t,0,2Pi},PlotStyle->Black];
Show[A,B]
RegionPlot3D[x^2+2*y^2<=z<=3-2*x^2-y^2&&x^2+y^2<=1,
{x,-1,1},{y,-1,1},{z,0,3},PlotPoints->80,Mesh->False,
PlotStyle->Orange]
\end{Verbatim}
\end{tcolorbox}

\subsection*{Una cuña circular}
Consideremos el sólido acotado por los gráficos de las ecuaciones
$r=3 \cos\theta,$ $z=-y$ y $z=0.$ Esta región se muestra a continuación

\begin{figure}[H]
	\centering
	\caption{Sólido determinado por $0\leq z \leq -y$, $r \leq 3 \cos \theta.$}
	\label{Sólido_determinado}
	\vspace{-4mm}
	\includegraphics[width=0.27\linewidth]{Fig/66.pdf}\hspace{1mm}
	% Eliminamos los vspace intermedios para evitar saltos de línea accidentales
	\vspace{-4mm}
	\includegraphics[width=0.27\linewidth]{Fig/67.pdf}\hspace{1mm}
	\vspace{-4mm}
	\includegraphics[width=0.27\linewidth]{Fig/68.pdf}
	
	\vspace{0.6cm} % <--- AJUSTA ESTE VALOR para bajar la fuente a tu gusto
	\parbox{0.95\linewidth}{\centering\fontsize{8pt}{10pt}\selectfont\textbf{Fuente:} elaboración propia.}
\end{figure}


La base del sólido es el semicírculo de radio $3$ que se parametriza en la forma $\textstyle\langle 3\cos\theta\cos\theta,\allowbreak\, 3\cos\theta\sin\theta \rangle$, y se simplifica como $\textstyle\langle 3\cos^2\theta,\allowbreak\, \tfrac{3}{2}\sin(2\theta) \rangle$, con $\theta \in [\pi, 2\pi]$. El interior de esta región se representa en el plano $xy$ como $\textstyle\langle 3u\cos^2\theta,\allowbreak\, \tfrac{3}{2}u\sin(2\theta) \rangle$, con $u \in [0,1]$ y $\theta \in [\pi, 2\pi]$. Un punto en la base del sólido tendrá la forma $\textstyle P = (3u\cos^2\theta,\allowbreak\, \tfrac{3}{2}u\sin(2\theta),\, 0)$, mientras que otro punto sobre el plano $z = -y$ tendrá la forma $\textstyle Q = (3u\cos^2\theta,\allowbreak\, \tfrac{3}{2}u\sin(2\theta),\allowbreak\, -\tfrac{3}{2}u\sin(2\theta))$. Formando una combinación convexa $\textstyle(1 - w)P + wQ$ entre estos dos puntos, se obtiene la función:
\vspace{-0.70mm}
{\small
\begin{equation*}
{\bf r}(\theta, u,w) = 
\bigg\langle 3u \cos^2\theta, \frac{3}{2}u \sin(2 \theta) , -\frac{3}{2}uw \sin(2 \theta)\bigg \rangle,
\end{equation*}
}
con $\theta \in  [\pi, 2\pi],$ $u\in [0,1]$ y $w\in [0,1],$ que parametriza el sólido.

\subsection*{Un sólido acotado por un cono y una esfera}
Las ecuaciones $z=\sqrt{x^2+y^2}$ y $x^2-4x+y^2+z^2=6$ 
representan un semicono y una esfera.
El eliminar $z$ de estas expresiones, se obtiene la ecuación
$\left( x-1\right) ^2+y^2=4$ que representa la curva de intersección de las superficies en el plano $xy$.
El interior de esta región se puede parametrizar en la forma
$x-1=2u\cos(t)$, $y=2u\sin(t)$, con $t \in [0,2\pi]$ y $u \in [0,1].$ 

\begin{figure}[H]
	\centering
	\caption{Semiesfera $z=\sqrt{6-x^2-y^2+4x}$ y semicono $z=\sqrt{x^2+y^2}$}
	\label{Semiesfera_y_semicono}
	\includegraphics[width=0.23\linewidth]{Img/74c.pdf}\hspace{3mm}
	\includegraphics[width=0.33\linewidth]{Fig/69.pdf} 
	\parbox{0.95\linewidth}{\centering\fontsize{9pt}{11pt}\selectfont{\bf Fuente:} elaboración propia.}
\end{figure}
\vskip0.2cm
Determinaremos la forma de los puntos sobre el semicono y sobre la esfera, posteriormente se hará una combinación convexa entre estos puntos con el fin de describir completamente el sólido. Sobre el semicono, se satisface que

\vspace{-7mm}
{\small
\[
 z=\sqrt{x^{2}+y^{2}}=\sqrt{\left( 1+2u\cos \left( t\right) \right)^{2}+\left( 2u\sin \left( t\right) \right) ^{2}}= \sqrt{1+4u\cos\left( t\right) +4u^{2}}.
\]
}

Entonces, los puntos sobre el cono (dentro del sólido) se determinan por

{\small
$$P\left( 1+2u\cos (t)  ,2u\sin (t) ,\sqrt{1+4u\cos\left( t\right) +4u^{2}} \right). $$
}

Reemplazando $x$ e $y$ en la ecuación de la esfera {\footnotesize $z=\sqrt{6 - x^{2} + 4x - y^{2}}$,} se obtiene {\footnotesize $z=\sqrt{6 - (1{+}2u\cos t)^2 + 4(1{+}2u\cos t) - (2u\sin t)^2} = \sqrt{9{+}4u\cos t{-}4u^2}$,} entonces los puntos de la forma {\footnotesize $Q(1{+}2u\cos t,\, 2u\sin t,\, \sqrt{9{+}4u\cos t{-}4u^2})$} yacen sobre la esfera. Por lo tanto, el sólido se describe al simplificar la expresión $(1{-}w)P{+}wQ$, para $w\in[0,1]$, de modo que la siguiente función describe el sólido, para $t\in[0,2\pi]$, $u\in[0,1]$ y $w\in[0,1]$. La figura \ref{region_conos} muestra esta situación.
\vspace{-04mm}
\begin{align*}
	\mathbf{r}(t,u,w) &= \left\langle 1{+}2u\cos t,\, 2u\sin t,\right. \\
	&\quad \left.(1{-}w)\sqrt{1{+}4u\cos t{+}4u^2} + w\sqrt{9{+}4u\cos t{-}4u^2} \right\rangle
\end{align*}


\vskip0.2cm
\noindent
\begin{minipage}{\textwidth}
	\centering	
	\captionof{figure}{\small Región determinada por $x^2+y^2 \leq z^2 \leq 6-x^2-y^2+4x$}
	\label{region_conos}
	\begin{minipage}{0.26\linewidth}
		\centering
		\includegraphics[width=\linewidth]{Img/74a.pdf}
	\end{minipage}
	\begin{minipage}{0.26\linewidth}
		\centering
		\includegraphics[width=\linewidth]{Img/74b.pdf}
	\end{minipage}
	
	\vspace{0.3em}
	\parbox{0.95\linewidth}{\centering\fontsize{8pt}{10pt}\selectfont\textbf{Fuente:} elaboración propia.}
\end{minipage}
\vskip0.2cm


La figura se obtiene ejecutando las siguientes instrucciones:



\begin{tcolorbox}[breakable,enhanced,title=\bf Instrucción \verb"Mathematica"] 
\begin{Verbatim}[formatcom=\letraverbatim]
ParametricPlot3D[{{1+2*u*Cos[t],2*u*Sin[t],
Sqrt[9+4*u*Cos[t]-4*u^2]},{1+2*u*Cos[t],
2*u*Sin[t],Sqrt[1+4*u*Cos[t]+4*u^2]}},
{u,0,1},{t,0,2*Pi},
PlotStyle->{RGBColor[0.8,0.3,0.3],Cyan},
MeshStyle->Gray,PlotPoints->60]
\end{Verbatim}
\end{tcolorbox}

\subsection*{Un sólido hueco}
Las expresiones $2x\leq x^{2}+y^{2}$ y $0\leq z\leq 9-x^{2}-y^{2},$
describen el sólido que se muestra a continuación.


\vskip0.2cm
\noindent
\begin{minipage}{\textwidth}
	\centering
	\captionof{figure}{\small Sólido interior al paraboloide y exterior al cilindro}
	\label{conos_naranjas}
	\begin{minipage}{0.22\linewidth}
		\centering
		\includegraphics[width=\linewidth]{Img/79a.pdf}
	\end{minipage}
	\begin{minipage}{0.22\linewidth}
		\centering
		\includegraphics[width=\linewidth]{Img/79c.pdf}
	\end{minipage}
	\begin{minipage}{0.22\linewidth}
		\centering
		\includegraphics[width=\linewidth]{Img/79b.pdf}
	\end{minipage}
	
	\vspace{0.3em}
	\parbox{0.95\linewidth}{\centering\fontsize{8pt}{10pt}\selectfont\textbf{Fuente:} elaboración propia.}
\end{minipage}
\vskip0.2cm


Este objeto consta de los puntos que yacen entre el plano $xy$ y el paraboloide $z=9-x^{2}-y^{2},$
exteriores al cilindro circular recto $2x=x^{2}+y^{2}.$
Esta última ecuación se reescribe como $\left( x-1\right) ^{2}+y^{2}=1$ y representa en el plano $xy$
una circunferencia. La figura \ref{conos_naranjas} se consigue ejecutando las siguientes instrucciones.
\vskip0.2cm

\begin{tcolorbox}[breakable,enhanced,title=\bf Instrucción \verb"Mathematica"]
\begin{Verbatim}[formatcom=\letraverbatim]
RegionPlot3D[x^2+y^2>2*x&&0<z<9-x^2-y^2,{x,-3,3},
{y,-3,3},{z,0,9},PlotPoints->90,Mesh->False,
PlotStyle->Orange]
\end{Verbatim}
\end{tcolorbox}
 

El paraboloide se proyecta en el plano $xy$ como una circunferencia con ecuación $x^{2}+y^{2}=9,$
la proyección del sólido en plano $xy$ contiene a los puntos interiores a una circunferencia y exteriores a la otra.


Para recorrer la circunferencia de radio $3$, usaremos {\small $2t$} como parámetro. La región entre las circunferencias en mención se parametriza
usando una combinación convexa entre dos puntos {\footnotesize $\left( 1+\cos (2t) ,\sin (2t)\right)$} y {\footnotesize $\left(3\cos(2t),3\sin(2t)\right),$} la cual se simplifica como 
\vspace{-3mm}
{\small
\begin{align*}
	{\bf r}(u,t)&=\big\langle u+(3-2 u)\cos (2t),\, (3-2u)\sin(2t)\big\rangle.
\end{align*}
}

\vskip0.2cm
\noindent
\begin{minipage}{\textwidth}
	\centering
	\captionof{figure}{${\bf r}(u,t)$}
	\label{fig:70}
	\begin{minipage}{0.24\linewidth}
		\centering
		\includegraphics[width=\linewidth]{Fig/70.pdf}
	\end{minipage}
	
	\vspace{0.3em}
	\parbox{0.95\linewidth}{\centering\fontsize{8pt}{10pt}\selectfont\textbf{Fuente:} elaboración propia.}
\end{minipage}
\vskip0.2cm


Las siguientes instrucciones muestran una animación el recorrido de las curvas.

\begin{tcolorbox}[breakable,enhanced,title=\bf Instrucción \verb"Mathematica"] 
\begin{Verbatim}[formatcom=\letraverbatim]
Manipulate[ParametricPlot[{{3*Cos[2*t],3*Sin[2*t]},
{1+Cos[2*t],Sin[2*t]}},{t,0,k},PlotRange->{-3,3}],
{k,Pi/20,Pi}]
\end{Verbatim}
\end{tcolorbox}

Con un punto $P(u+(3-2u)\cos(2t),(3-2u)\sin(2t),0)$ en la base del sólido 
se obtiene la coordenada $z$ de un punto $Q$ sobre el paraboloide,
\begin{eqnarray*}
z &=& 9-(u+(3-2u)\cos(2t))^2-((3-2u)\sin(2t))^{2} \\
&=& u (12-5 u+2(2u-3)\cos(2t)).
\end{eqnarray*}
El sólido se parametriza formando segmentos rectilíneos que unen a $P$
con $Q,$ esto se logra mediante la combinación convexa $(1-v)P+vQ,$
en donde $Q$ tiene la forma
$Q(u+(3-2u)\cos(2t),(3-2u)\sin(2t),u (12-5 u+2(2u-3)\cos(2t))).$ Esto nos lleva a definir la función 
\begin{multline} \label{solidohueco}
{\bf r}(t,u,v) = \big \langle u+(3-2u)\cos(2t),(3-2u)\sin(2t), \\ uv(12-5 u+2(2u-3)\cos(2t)) \big \rangle
\end{multline}
En donde $t\in [0,\pi]$, $u\in [0,1]$ y $v \in [0,1]$.

\subsection*{Un sólido hueco II} \label{solidohueco3}

Las ecuaciones $z=4-y^{2},$ $z=4-x,$ $z=0,$ $x=0,$ representan un cilindro parabólico y 
tres planos. Estas superficies acotan un sólido $E,$
que se perfora por el cilindro $\left( z-2\right) ^{2}+y^{2}=1,$ obteniendo la región 
que se muestra en la figura.


\vskip0.2cm
\noindent
\begin{minipage}{\textwidth}
	\centering
	\captionof{figure}{$z=4-y^2$, $z=4-x$}
	\label{fig:71}
	\begin{minipage}{0.30\linewidth}
		\centering
		\includegraphics[width=\linewidth]{Fig/71.pdf}
	\end{minipage}
	
	\vspace{0.3em}
	\parbox{0.95\linewidth}{\centering\fontsize{8pt}{10pt}\selectfont\textbf{Fuente:} elaboración propia.}
\end{minipage}
\vskip0.2cm

 
%\vspace{-4mm}
La figura del sólido se obtiene ejecutando las siguientes instrucciones.

\begin{tcolorbox}[breakable,enhanced,title=\bf Instrucción \verb"Mathematica"] 
\begin{Verbatim}[formatcom=\letraverbatim]
RegionPlot3D[0<x<4-z&&(z-2)^2+y^2>1&&0<z<4-y^2&&
-2<y<2,{x,0,4},{y,-2,2},{z,0,4},PlotPoints->90,
PlotStyle->Orange,Mesh->False]
\end{Verbatim}
\end{tcolorbox}

\noindent
\begin{minipage}{\textwidth}
	\centering
	\captionof{figure}{\small Región determinada por $0\leq z \leq 4-x$, $z\leq 4-y^2,$ $1\leq (z-2)^2+y^2.$}
	\begin{minipage}{0.22\linewidth}
		\centering
		\includegraphics[width=\linewidth]{Img/5a.pdf}
	\end{minipage}
	\hspace{1em}
	\begin{minipage}{0.22\linewidth}
		\centering
		\includegraphics[width=\linewidth]{Img/5c.pdf}
	\end{minipage}
	\hspace{1em}
	\begin{minipage}{0.22\linewidth}
		\centering
		\includegraphics[width=\linewidth]{Img/5b.pdf}
	\end{minipage}
	
	\vspace{0.3em}
	\parbox{0.95\linewidth}{\centering\fontsize{8pt}{10pt}\selectfont\textbf{Fuente:} elaboración propia.}
\end{minipage}


La proyección del sólido en el plano $yz$ corresponde a una circunferencia y una porción de 
parábola, que se representan en la forma $\ve{r}_1(t)=\la 0,\cos (t),2+\sin (t)\ra,$ con 
$t\in [0,2\pi]$ y $\ve{r}_2(t)=\la 0,t,4-t^{2}\ra,$ con $t\in \lbrack -2,2],$ respectivamente.
Inicialmente parametrizaremos la base del sólido en el plano $yz$. Esta región a su vez 
la dividiremos en dos porciones.

\begin{figure}[H] 
	\centering
	\caption{Región determinada por $0\leq z \leq 4-y^2$, $1\leq (z-2)^2+y^2.$}
	\label{Región_determinada}
	\begin{minipage}[t]{0.35\linewidth} 
		\centering
		\includegraphics[width=\linewidth]{Fig/72.pdf}
	\end{minipage}
	\hspace{1em}
	\begin{minipage}[t]{0.35\linewidth} 
		\centering
		\includegraphics[width=\linewidth]{Fig/73.pdf}
	\end{minipage}
	\parbox{0.95\linewidth}{\centering\fontsize{9pt}{11pt}\selectfont{\bf Fuente:} elaboración propia.}
\end{figure}
%\vskip0.1cm
Para describir las regiones anteriores unificaremos el dominio de las parametrizaciones dadas y luego construir segmentos rectilíneos entre las curvas.
\vskip2cm
\vspace{-10mm}

\noindent
\textcolor{black}{\scriptsize\ding{110}}\quad {\bf Región 1.} Para aplicar el intervalo $[-2,2]$ en $[0,\pi ],$ se utiliza la relación
$v=\frac{\pi }{4}\left(u+2\right)$; además, con el intercambio de $t$ por $\pi -t$ se invierte el recorrido. 

\vskip0.1cm
El cambio de $t$ por $2-\frac{4t}{\pi}$ en la parametrización de la parábola,
unifica el dominio de las funciones. Entonces, con una combinación convexa de los puntos 
$P(0,\cos (t),2+\sin(t))$ y $Q(0,2-\frac{4t}{\pi},\frac{16t}{\pi}\left(1-\frac{t}{\pi}\right))$, se obtiene la función que parametriza la región en mención, esto es
\begin{multline*}
{\bf r}_1(t,u)  =\big \langle 0, (1-u)\cos (t) +u\left(2-\frac{4t}{\pi}\right),\\ \frac{16(\pi-t)tu}{\pi^2}-(u-1)(\sin(t)+2) \big\rangle, 
\end{multline*} 
para $t\in [0,\pi ]$ y $u\in [0,1]$.

\vspace{2mm}

\noindent
\textcolor{black}{\scriptsize\ding{110}}\quad {\bf Región 2.} El segmento que une los puntos $\left(0,-2,0 \right) $ con $\left(0,2,0 \right),$
se parametriza como $\langle 0,4(t-\pi)/\pi-2,0,0 \rangle $,
con $t\in [\pi ,2\pi ].$ Trazando segmentos rectilíneos desde el segmento anterior 
hasta la semicircunferencia, se consigue parametrizar el conjunto de interés; esto se logra 
formando una combinación convexa entre los puntos
$(0,\cos (t),2+\sin (t))$ y $(0,4(t-\pi)/\pi-2,0)$ por esta razón, la función 
$${\bf r}_2\left( t,u\right)  = \left\langle 0,\left(
1-u\right) \cos (t) +u\left( \frac{4\left( t-\pi \right) }{\pi }%
-2\right) ,\left( 1-u\right) \left( 2+\sin (t) \right)
\right\rangle ,$$ con $t\in \lbrack \pi ,2\pi ]$ y $u\in \lbrack 0,1],$ parametriza la región considerada.


Con las funciones ${\bf r}_1$ y ${\bf r}_2$ se representará la región sobre el 
plano $z=4-x,$ esto cambia solamente la primera función componente de estas funciones. 
\begin{multline*}
{\bf r}_3(u,t)=\bigg \langle 4-\frac{16 (\pi-t) t u}{\pi ^2}+(u-1) (\sin (t)+2),
(1-u)\cos (t) \\ +u\left(2-\frac{4t}{\pi}\right), \frac{16 (\pi -t) t u}{\pi ^2}-(u-1) (\sin (t)+2) \bigg\rangle, t\in[0,\pi],\, u\in[0,1].
\end{multline*}  \begin{multline*}
{\bf r}_4(u,t)= \big \langle 4-(1-u)(\sin(t)+2),\left(\frac{4(t-\pi)}{\pi}-2\right)u+(1-u)\cos(t),\\
(1-u)(\sin (t)+2)\big \rangle,\,t\in[\pi,2\pi],\, u\in[0,1].
\end{multline*}

Con las siguientes instrucciones se obtiene el gráfico de las superficies que acotan el sólido.


%\fontsize{9.5}{10}\selectfont
\begin{tcolorbox}[breakable,enhanced,title=\bf Instrucción \verb"Mathematica"] 
\begin{Verbatim}[formatcom=\letraverbatim]
y1=(1-u)*Cos[t]+u*(2-4*t/Pi);
z1=16*(Pi-t)*u*t/Pi^2-(u-1)*(2+Sin[t]);
{y2,z2}={(1-u)*Cos[t]+u*(4*(t-Pi)/Pi-2),
(1-u)*(2+Sin[t])};
A1=ParametricPlot3D[{{0,y1,z1},{4-z1,y1,z1}},{t,0,Pi},
{u,0,1},Mesh->False,PlotStyle->Orange];
A2=ParametricPlot3D[{{0,y2,z2},{4-z2,y2,z2}},
{t,Pi,2*Pi},{u,0,1},Mesh->False,PlotStyle->Cyan];
r5=(1-v)*{0,y1,z1}+v*{4-z1,y1,z1}/.u->1;
A3=ParametricPlot3D[r5,{t,0,Pi},{v,0,1},
PlotStyle->Orange,Mesh->False];
Show[A1,A2,A3]
\end{Verbatim}
\end{tcolorbox}

El sólido quedará parametrizado al hacer una combinación convexa entre dos puntos: 
uno en la base del sólido (del plano $yz$) y otro punto sobre el plano $z=4-x$. 
Es por lo expuesto antes, que una parametrización del sólido es:
\begin{equation} \label{solidohueco2}
{\bf r}(u,t,w) = \left \{ \begin{array}{cc}
(1-w){\bf r}_1(u,t)+w {\bf r}_3(u,t) & t \in [0,\pi] \\[0.2cm]
(1-w){\bf r}_2(u,t)+w {\bf r}_4(u,t) & t \in [\pi,2\pi] 
\end{array}\right . 
\end{equation}

\subsection*{Sólidos con una cardioide en su base} \index{Cardioide}
\noindent
\textcolor{black}{\scriptsize\ding{110}}\quad La cardioide con ecuación $r=1+\cos (t) $
se parametriza en el plano $xy$ como  $\left\langle (1+\cos(t)) \cos(t),(1+\cos (t))\sin (t) \right\rangle ,$ 
$t\in [0,2\pi];$ para representar su interior se hace el siguiente cambio sutil
$\big\langle u(1+\cos(t)) \cos(t),u(1+\cos (t))\sin (t) \big\rangle ,$  con $u\in[0,1].$
Estos puntos se unen a $(1,0,1)$ utilizando el par de funciones inyectivas en 
$[0,1],$ $f(v)=\sqrt[5]{(1-v) ^{2}}$ y $g(v) =v^{2},$ lo cual parametriza el sólido que se muestra a continuación.

\begin{figure}[H]
\begin{center}
	\captionof{figure}{\small Gráfico de ${\bf r}\left(t,u,w\right)$}
\includegraphics[height=4.5cm]{Img/76b.pdf} 
\includegraphics[height=4.5cm]{Img/76a.pdf}
\includegraphics[height=4.5cm]{Img/76c.pdf}
\end{center}
\parbox{0.95\linewidth}{\centering\fontsize{9pt}{11pt}\selectfont{\bf Fuente:} elaboración propia.}
\end{figure}

La combinación que genera el sólido es:
$${\bf r}(t,u,w) =\sqrt[5]{(1-v)^2}\, u(1+\cos(t)) (\cos(t),\sin(t),0) +v^{2}(1,0,1),$$
con $t \in [0,2\pi]$, $u \in [0,1]$ y $w \in [0,1].$ 
Las siguientes instrucciones, con $u=1$, permiten obtener el gráfico anterior.

\begin{tcolorbox}[breakable,enhanced,title=\bf Instrucción \verb"Mathematica"] 
\begin{Verbatim}[formatcom=\letraverbatim]
r=(1-v)^(2/5)*{(1+Cos[t])*Cos[t],(1+Cos[t])*Sin[t],0}
+v^2*{1,0,1};
p={r[[1]],r[[2]],0};
ParametricPlot3D[{r,p},{v,0,1},{t,0,2Pi},PlotPoints
->80,PlotStyle->{Orange,Cyan},MeshStyle->Gray]
\end{Verbatim}
\end{tcolorbox}

\vspace{2mm}

\noindent
\textcolor{black}{\scriptsize\ding{110}}\quad De manera similar al ejemplo anterior, se puede utilizar
la función 
$${\bf r}(t,u,v)=\sqrt[4]{(1-v)^3}\, u(1+\cos(t)) (\cos(t),\sin(t),0)+\sqrt[4]{v^3}(1,0,1),$$
con $t \in [0,2\pi]$, $u \in [0,1]$ y $w \in [0,1].$ Para $u=1$ se genera la superficie cuyo gráfico se muestra
a continuación.

\begin{figure}[H]
\begin{center}
	\captionof{figure}{\small Gráfico de ${\bf r}\left(t,u,w\right)$}
\includegraphics[height=4.5cm]{Img/77a.pdf} 
\includegraphics[height=4.5cm]{Img/77b.pdf}
\end{center}
\parbox{0.95\linewidth}{\centering\fontsize{9pt}{10pt}\selectfont{\bf Fuente:} elaboración propia.}
\end{figure}

Con las siguientes instrucciones se puede obtener el gráfico.

\begin{tcolorbox}[breakable,enhanced,title=\bf Instrucción \verb"Mathematica"]
\begin{Verbatim}[formatcom=\letraverbatim]
r=(1-v)^(3/4)*{(1+Cos[t])*Cos[t],(1+Cos[t])*Sin[t],0}
+v^(3/4)*{1,0,1};
p={r[[1]],r[[2]],0};
ParametricPlot3D[{r,p},{v,0,1},{t,0,2Pi},PlotPoints
->80,PlotStyle->{Orange,Cyan},MeshStyle->Gray]
\end{Verbatim}
\end{tcolorbox}
  
\subsection*{Cono cardioidico} \index{Cono cardioidico}
El interior de la cardioide con ecuación polar $r=1+\cos(t),$ \index{Cardioide}
se puede parametrizar en el plano $xy$, como 
$\left\langle u\left(1+\cos (t) \right) \cos (t) ,u\left( 1+\cos \left(
t\right) \right) \sin (t) ,0\right\rangle $. El sólido
que se genera mediante una combinación convexa de los puntos de esta región 
con el punto $\left( 1,0,1\right) $ es un \emph{cono cardioidico}\footnote{Este término fué acuñado por el profesor Juan Zambrano, de la Universidad Distrital Francisco José de Caldas.}.

\begin{figure}[H]
\begin{center}
	\captionof{figure}{\small Cono cardioidico}
\includegraphics[height=4.5cm]{Img/75a.pdf} 
\includegraphics[height=4.5cm]{Img/75b.pdf}
\includegraphics[height=4.5cm]{Img/75c.pdf}
\end{center}
\parbox{0.95\linewidth}{\centering\fontsize{9pt}{10pt}\selectfont\textbf{Fuente:} elaboración propia.}
\end{figure}

Es por esto que al simplificar la combinación 
$$\left( 1-w\right) \left( u\left( 1+\cos (t) \right) \cos \left(
t\right) ,u\left( 1+\cos (t) \right) \sin (t),0\right) +w\left( 1,0,1\right) $$
se obtiene la función 
${\bf r}(t,u,w)=\big\langle u\cos (t) \left(
1+\cos (t) \right) \left( 1-w\right) +w,\linebreak (1-w)u\left( 1+\cos (t) \right) \sin (t),w \big\rangle ,$
que describe el mencionado sólido, para $t\in [0,2\pi]$, $u\in [0,1]$ y $w\in [0,1].$
Los gráficos se obtienen con ayuda de las siguientes instrucciones.

\begin{tcolorbox}[breakable,enhanced,title=\bf Instrucción \verb"Mathematica"]
\begin{Verbatim}[formatcom=\letraverbatim]
p={(1+Cos[t])*Cos[t],(1+Cos[t])*Sin[t],0};
r=(1-u)*p+u*{1,0,1};
ParametricPlot3D[{r,p*u},{u,0,1},{t,0,2Pi},
Ticks->{{0,1,2},{-1,0,1},{0,1}},PlotStyle
->{Orange,Cyan},MeshStyle->Gray,PlotPoints->80]
\end{Verbatim}
\end{tcolorbox}

%\clearpage
\begin{ejer}
Parametrice y grafique en \verb"Mathematica" las siguientes regiones.
\end{ejer}

\begin{enumerate}
\item La región del plano comprendida entre las elipses $\frac{x^2}{16}+\frac{y^2}{9}=1$ y $(x-1)^2+\frac{y^2}{4}=1$.

\vskip0.2cm
\vspace{-6mm}
\begin{figure}[H]
	\centering
	\caption{Región acotada por dos elipses}
	\label{Región_acotada_dos_elipses}
	\includegraphics[width=0.4\linewidth]{Fig/74.pdf}
	\parbox{0.95\linewidth}{\centering\fontsize{8pt}{10pt}\selectfont\textbf{Fuente:} elaboración propia.}
\end{figure}

\item La región del plano acotada por las elipses $\frac{x^2}{16}+\frac{y^2}{9}=1$ y $(x-1)^2+\frac{y^2}{4}=1$,
se proyecta sobre el paraboloide hiperbólico $3z=y^2-x^2-2y$.

\begin{figure}[H]
	\centering
	\caption{Superficie sobre el paraboloide hiperbólico $3z = y^2 - x^2 - 2y$.}
	\label{Superficie_paraboloide_hiperbólico}
	\begin{minipage}[t]{0.30\linewidth}
		\centering
		\includegraphics[width=\linewidth]{Fig/75.pdf}
	\end{minipage}%
	\hspace{1em}
	\begin{minipage}[t]{0.30\linewidth}
		\centering
		\includegraphics[width=\linewidth]{Fig/76.pdf}
	\end{minipage}
	\parbox{0.95\linewidth}{\centering\fontsize{8pt}{10pt}\selectfont\textbf{Fuente:} elaboración propia.}
\end{figure}

\item Para $a_1=\cosh 2$, $b_1=\sinh(2)$, $a_2=\cos(1/2)$ y $b_2=\sin(1/2)$, la región del plano interior a la elipse
$\frac{x^2}{a_1^2}+\frac{y^2}{b_1^2}= 1$ y a la hipérbola $\frac{x^2}{a_2^2}-\frac{y^2}{b_2^2}= 1.$
\vspace{-1mm}
\begin{figure}[H]
	\centering
	\caption{Región acotada por una elipse y una hipérbola}
	\label{Región_acotada_elipse_hipérbola}
	\includegraphics[width=0.3\linewidth]{Fig/77.pdf}
	\parbox{0.95\linewidth}{\centering\fontsize{8pt}{10pt}\selectfont\textbf{Fuente:} elaboración propia.}
\end{figure}
\vspace{-2mm}
{\bf Sugerencia:} Utilice coordenadas curvilíneas
\vskip0.2cm
\index{Coordenadas curvilíneas}
\item Para $v=\pm \ln(5)$ parametrizar, con una sola función la región exterior a las dos circunferencias
con ecuaciones $(x-\coth v)^2+y^2=\frac{24}{144}$.

\begin{figure}[H]
	\begin{center}
		\captionof{figure}{\small Región exterior a dos circunferencias.}
		\includegraphics[height=4.5cm]{Img/T_71a.pdf}
	\end{center}
	\parbox{0.95\linewidth}{\centering\fontsize{8pt}{10pt}\selectfont\textbf{Fuente:} elaboración propia.}
\end{figure}
\vspace{-5mm}
{\bf Sugerencia:} Utilice coordenadas curvilíneas.

\item Para $v=\pm \ln(5)$ parametrizar, con una sola función, la región sobre el paraboloide hiperbólico 
$6z=x^2-y^2$, exterior a los cilindros con ecuaciones $(x-\coth v)^2+y^2=\frac{24}{144}$.
\begin{figure}[H]
\begin{center}
	\captionof{figure}{\small Región sobre el paraboloide hiperbólico.}
\includegraphics[height=4.5cm]{Img/T_71b.pdf} \hskip0.5cm
\includegraphics[height=4.5cm]{Img/T_71c.pdf}
\end{center}
\parbox{0.95\linewidth}{\centering\fontsize{8pt}{10pt}\selectfont\textbf{Fuente:} elaboración propia.}

\end{figure} {\bf Sugerencia:} Utilice coordenadas curvilíneas
\begin{comment}
Clear[x,y,z,u,v];
x=Sinh[v]/(Cosh[v]-Cos[u]);
y=Sin[u]/(Cosh[v]-Cos[u]);
z=(y^2+6x-x^2)/6;
ParametricPlot3D[{x,y,z},{u,0,2*Pi},{v,-Log[3],Log[3]},
PlotStyle->Orange,Mesh->20,PlotPoints->95,
MeshStyle->Gray,PlotRange->{{-4,10},{-4,4},{-4,4.2}},BoxRatios->{1,1,1/2}]
\end{comment}

\vskip0.2cm
\item El sólido acotado por el cilindro elíptico $2x^2+y^2=16$ y los parabo- \\ loides hiperbólicos
$x^2-y^2= 4z+12$ y $x^2-2y^2=4z-20$. 

%\vspace{-4mm}
\vskip0.2cm
\noindent
\begin{minipage}{\textwidth}
	\centering
	\captionof{figure}{\small $2x^2+y^2\leq 16$, $ x^2-y^2-12 \leq 4z\leq 20 + x^2-2y^2.$}
	\begin{minipage}{0.17\linewidth}
		\centering
		\includegraphics[width=\linewidth]{Img/T_56a.pdf}
	\end{minipage}
	\hspace{1em}
	\begin{minipage}{0.17\linewidth}
		\centering
		\includegraphics[width=\linewidth]{Img/T_56b.pdf}
	\end{minipage}
	
	\vspace{0.3em}
	\parbox{0.95\linewidth}{\centering\fontsize{8pt}{10pt}\selectfont\textbf{Fuente:} elaboración propia.}
\end{minipage}
\vskip0.2cm
%RegionPlot3D[2*x^2+y^2<=16&&x^2-y^2-12<=4*z<=20+x^2-2*y^2,{x,-3,3},{y,-4,4},{z,-7,7},PlotPoints->90,Mesh->False]
\item El sólido acotado por las figuras del plano $5z+x+y=20$ y el semicono 
$z=\sqrt{x^2+y^2}.$ 


\vskip0.2cm
\noindent
\begin{minipage}{\textwidth}
	\centering
	\captionof{figure}{\small $\sqrt{x^2+y^2} \leq z \leq (20-x-y)/5$}
	\begin{minipage}{0.22\linewidth}
		\centering
		\includegraphics[width=\linewidth]{Img/T_39a.pdf}
	\end{minipage}
	
	\vspace{0.3em}
	\parbox{0.95\linewidth}{\centering\fontsize{8pt}{10pt}\selectfont\textbf{Fuente:} elaboración propia.}
\end{minipage}
\vskip0.2cm



%\begin{comment}
%RegionPlot3D[z>=Sqrt[x^2+y^2]&&5*z<20-y-x,{x,-6,4},{y,-6,4},{z,0,6},PlotPoints->90,Mesh->False]
%\end{comment}

\item La porción del cilindro elíptico con ecuación $4x^2+9y^2=36$ acotado por los planos 
$3z+x=9$ y $3z-y=-9$.

\vskip0.2cm
\noindent
\begin{minipage}{\textwidth}
	\centering
	\captionof{figure}{\small $4x^2+9y^2=36,$ $-9+y \leq 3z \leq 9-x.$}
	\begin{minipage}{0.14\linewidth}
		\centering
		\includegraphics[width=\linewidth]{Img/T_38a.pdf}
	\end{minipage}
	
	\vspace{0.3em}
	\parbox{0.95\linewidth}{\centering\fontsize{8pt}{10pt}\selectfont\textbf{Fuente:} elaboración propia.}
\end{minipage}
\vskip0.2cm




%\begin{comment}
%\begin{Verbatim}
%a=ParametricPlot3D[{{3*Cos[t],2*Sin[t],3-Cos[t]},{3*Cos[t],2*Sin[t],-3+Sin[t]}},
%{t,0,2*Pi},PlotStyle->{Directive[Black,Thickness->0.015],Directive[Black,Thickness->0.015]}];
%b=ParametricPlot3D[(1-u)*{3*Cos[t],2*Sin[t],3-Cos[t]}+u{3*Cos[t],2*Sin[t],-3+Sin[t]},
%{t,0,2*Pi},{u,0,1},MeshStyle->Gray];
%Show[a,b]
%\end{Verbatim}
%\end{comment}
\end{enumerate}
